% Generated by tools/build.py -- do not edit; edit data/toc-compilers/*.txt
% No preamble here: this file is \input by main.tex.
\chapter{Unit 8: Theory of Computation and Compilers}

\section{Formal Languages, Diagonalization and Finite Automata (DFA, NFA, minimization, Moore and Mealy machines)}

\begin{question}{H}
The minimum number of states in a DFA over \{0, 1\} accepting strings in which the number of 0s is divisible by 3 is
\optsshort{3}{2}{4}{6}
\end{question}

\begin{question}{H}
The minimum number of states in a DFA that accepts binary numbers divisible by 5 (reading the most significant bit first, and treating the empty string as 0) is
\optsshort{5}{4}{6}{10}
\end{question}

\begin{question}{H}
The minimum number of states in a DFA over \{0, 1\} accepting all strings that end with 01 is
\optsshort{2}{5}{4}{3}
\end{question}

\begin{question}{H}
The minimum number of states in a DFA over \{0, 1\} accepting all strings that contain 101 as a substring is
\optsshort{3}{6}{5}{4}
\end{question}

\begin{question}{H}
The minimum number of states in a DFA over \{0, 1\} accepting strings whose third symbol from the right end is 1 is
\optsshort{4}{8}{6}{3}
\end{question}

\begin{question}{H}
A minimal DFA for strings over \{0, 1\} in which the nth symbol from the right is 1 has
\optsshort{n\textsuperscript{2} states}{2n states}{n + 1 states}{2\textsuperscript{n} states}
\end{question}

\begin{question}{H}
The minimum number of states in a DFA over \{0, 1\} accepting strings with an even number of 0s and an odd number of 1s is
\optsshort{2}{3}{6}{4}
\end{question}

\begin{question}{H}
The minimum number of states in a DFA over \{0, 1\} accepting strings that start with 1 and end with 0 is
\optsshort{5}{4}{3}{6}
\end{question}

\begin{question}{M}
The minimum number of states in a DFA over \{0, 1\} accepting all strings of length exactly 3 is
\optsshort{3}{4}{8}{5}
\end{question}

\begin{question}{M}
The minimum number of states in a DFA over \{0, 1\} accepting strings that contain exactly two 1s is
\optsshort{6}{3}{4}{5}
\end{question}

\begin{question}{M}
The minimum number of states in a DFA over \{a, b\} accepting strings in which the number of a's is even and the number of b's is divisible by 3 is
\optsshort{8}{12}{6}{5}
\end{question}

\begin{question}{M}
The minimum number of states in a DFA over \{0, 1\} accepting strings in which the first and last symbols are the same (strings of length at least 1) is
\optsshort{4}{3}{5}{6}
\end{question}

\begin{question}{H}
In the subset construction, an NFA with n states can result in a DFA with at most
\optsshort{n states}{2\textsuperscript{n} states}{n\textsuperscript{2} states}{2n states}
\end{question}

\begin{question}{H}
Given below are two statements.\\
Statement I: An NFA with $\varepsilon$-transitions accepts a larger class of languages than an NFA without them.\\
Statement II: Every NFA can be converted into an equivalent DFA.
\optslong{Both Statement I and Statement II are true}{Both Statement I and Statement II are false}{Statement I is true but Statement II is false}{Statement I is false but Statement II is true}
\end{question}

\begin{question}{H}
Match List I with List II.\\
List I (Machine): A. Moore machine  B. Mealy machine  C. DFA  D. Two-way DFA\\
List II (Property): I. Accepts only regular languages despite moving the head in both directions  II. Output depends on the current state and the current input  III. Output depends only on the current state  IV. Has no output other than accept or reject
\optsshort{A-III, B-II, C-IV, D-I}{A-III, B-IV, C-II, D-I}{A-III, B-II, C-I, D-IV}{A-II, B-III, C-IV, D-I}
\end{question}

\begin{question}{M}
For an input string of length n, a Moore machine produces an output string of length
\optsshort{n + 1}{n}{n $-$ 1}{2n}
\end{question}

\begin{question}{M}
Given below are two statements.\\
Statement I: Every Mealy machine can be converted into an equivalent Moore machine, possibly with more states.\\
Statement II: Every Moore machine can be converted into an equivalent Mealy machine with the same number of states.
\optslong{Both Statement I and Statement II are true}{Both Statement I and Statement II are false}{Statement I is true but Statement II is false}{Statement I is false but Statement II is true}
\end{question}

\begin{question}{H}
In a DFA, the transition function $\delta$ maps
\optsshort{Q $\times$ ($\Sigma$ $\cup$ \{$\varepsilon$\}) to 2\textsuperscript{Q}}{Q $\times$ $\Sigma$ to Q}{Q to $\Sigma$}{Q $\times$ $\Sigma$ to 2\textsuperscript{Q}}
\end{question}

\begin{question}{H}
The Myhill-Nerode theorem implies that a language L is regular if and only if
\optslong{L has a regular grammar with $\varepsilon$-productions}{L is finite}{the relation "xz $\in$ L exactly when yz $\in$ L for all z" has finitely many equivalence classes}{L is accepted by a PDA}
\end{question}

\begin{question}{M}
Given below are two statements: one is labelled as Assertion A and the other is labelled as Reason R.\\
Assertion A: The minimal DFA for a regular language is unique up to renaming of states.\\
Reason R: Its states correspond one-to-one to the Myhill-Nerode equivalence classes of the language.
\optslong{Both A and R are true and R is the correct explanation of A}{Both A and R are true but R is NOT the correct explanation of A}{A is true but R is false}{A is false but R is true}
\end{question}

\begin{question}{M}
In DFA minimization by the table-filling (partitioning) method, two states are distinguishable if
\optslong{both are accepting}{some string leads one of them to an accepting state and the other to a non-accepting state}{they are reachable from the start state}{they have the same number of transitions}
\end{question}

\begin{question}{H}
The set of all languages over the alphabet \{0, 1\} is
\optsshort{empty}{uncountable}{finite}{countably infinite}
\end{question}

\begin{question}{H}
Given below are two statements: one is labelled as Assertion A and the other is labelled as Reason R.\\
Assertion A: There exist languages that are not accepted by any Turing machine.\\
Reason R: A Turing machine can simulate any finite automaton.
\optslong{Both A and R are true and R is the correct explanation of A}{Both A and R are true but R is NOT the correct explanation of A}{A is true but R is false}{A is false but R is true}
\end{question}

\begin{question}{M}
Cantor's diagonal argument is used to prove that
\optslong{the set of real numbers in [0, 1] is uncountable}{every TM halts}{every language is regular}{the set of integers is countable}
\end{question}

\begin{question}{M}
Russell's paradox concerns
\optslong{the number of states in a DFA}{the pumping lemma}{the set of all sets that do not contain themselves}{the halting of Turing machines}
\end{question}

\begin{question}{M}
Given below are two statements.\\
Statement I: The set of all subsets of the natural numbers is countable.\\
Statement II: The set of all strings over a finite alphabet is countable.
\optslong{Both Statement I and Statement II are true}{Both Statement I and Statement II are false}{Statement I is true but Statement II is false}{Statement I is false but Statement II is true}
\end{question}

\begin{question}{M}
The number of strings of length at most n over the alphabet \{0, 1\}, including the empty string, is
\optsshort{2\textsuperscript{n}}{n\textsuperscript{2}}{2\textsuperscript{n+1} $-$ 1}{2\textsuperscript{n} $-$ 1}
\end{question}

\begin{question}{M}
The number of strings of length 5 over \{0, 1\} that do not contain two consecutive 1s is
\optsshort{13}{8}{21}{16}
\end{question}

\begin{question}{M}
A DFA with n states accepts a non-empty language if and only if it accepts some string of length
\optsshort{n\textsuperscript{2}}{less than n}{greater than 2n}{exactly n}
\end{question}

\begin{question}{M}
Given below are two statements.\\
Statement I: A DFA with n states accepts an infinite language if and only if it accepts some string of length between n and 2n $-$ 1.\\
Statement II: Whether a DFA accepts an infinite language is undecidable.
\optslong{Both Statement I and Statement II are true}{Both Statement I and Statement II are false}{Statement I is true but Statement II is false}{Statement I is false but Statement II is true}
\end{question}

\begin{question}{M}
The minimum number of states of a DFA over \{0, 1\} that accepts strings of length divisible by 2 or by 3 is
\optsshort{7}{6}{12}{5}
\end{question}

\begin{question}{M}
Which of the following is true of a dead (trap) state in a DFA?
\optslong{It is the start state}{It is always accepting}{Every transition from it leads back to itself, and it is non-accepting}{It has no outgoing transitions}
\end{question}

\begin{question}{M}
Lexical analyzers are typically implemented using
\optsshort{pushdown automata}{linear-bounded automata}{finite automata}{Turing machines}
\end{question}

\begin{question}{L}
The language accepted by a finite automaton is called a
\optsshort{context-free language}{recursive language}{regular language}{context-sensitive language}
\end{question}

\begin{question}{M}
Given below are two statements: one is labelled as Assertion A and the other is labelled as Reason R.\\
Assertion A: Finite automata cannot recognise the language \{0\textsuperscript{n} 1\textsuperscript{n} : n $\geq$ 0\}.\\
Reason R: A finite automaton has no way to remember an unbounded count with finitely many states.
\optslong{Both A and R are true and R is the correct explanation of A}{Both A and R are true but R is NOT the correct explanation of A}{A is true but R is false}{A is false but R is true}
\end{question}

\begin{question}{M}
Given below are two statements.\\
Statement I: An NFA accepts a string if every possible computation path ends in an accepting state.\\
Statement II: An NFA may have several transitions on the same symbol from one state.
\optslong{Both Statement I and Statement II are true}{Both Statement I and Statement II are false}{Statement I is true but Statement II is false}{Statement I is false but Statement II is true}
\end{question}

\begin{question}{M}
The minimum number of states in a DFA over \{0, 1\} accepting strings that do not contain two consecutive 1s is
\optsshort{4}{2}{5}{3}
\end{question}

\begin{question}{M}
Given below are two statements.\\
Statement I: The complement of a DFA's language is obtained by reversing all its transitions.\\
Statement II: The complement of an NFA's language is always obtained by swapping its accepting and non-accepting states.
\optslong{Both Statement I and Statement II are true}{Both Statement I and Statement II are false}{Statement I is true but Statement II is false}{Statement I is false but Statement II is true}
\end{question}

\begin{answers}
\ans{1}{1}{The DFA counts 0s modulo 3; 1s do not change the state.}
\ans{2}{1}{One state for each remainder 0--4; reading bit b takes remainder r to (2r + b) mod 5.}
\ans{3}{4}{States track the longest suffix that is a prefix of 01: $\varepsilon$, 0 and 01.}
\ans{4}{4}{States for $\varepsilon$, 1, 10 and 101 (the last one is a trap accepting state).}
\ans{5}{2}{The DFA must remember the last three symbols: 2\textsuperscript{3} = 8 states. An NFA needs only 4.}
\ans{6}{4}{This family shows that the subset construction can cause an exponential blow-up.}
\ans{7}{4}{One state for each combination of parities.}
\ans{8}{2}{A start state, a dead state (string begins with 0), and two states for "starts with 1, last symbol 0 or 1".}
\ans{9}{4}{States count the length 0, 1, 2, 3, plus a dead state for longer strings.}
\ans{10}{3}{States for zero, one and two 1s, plus a dead state for more than two.}
\ans{11}{3}{Product of a mod-2 counter and a mod-3 counter: 2 $\times$ 3 = 6.}
\ans{12}{3}{A start state, then for each first symbol a state for "last symbol same" and one for "last symbol different".}
\ans{13}{2}{Each DFA state is a subset of the NFA states.}
\ans{14}{4}{DFAs, NFAs and $\varepsilon$-NFAs all accept exactly the regular languages.}
\ans{15}{1}{Moore and Mealy machines are finite-state transducers.}
\ans{16}{1}{It produces an output in the start state before reading any input; a Mealy machine produces n outputs.}
\ans{17}{1}{The equivalence ignores the extra initial output of the Moore machine.}
\ans{18}{2}{For an NFA, $\delta$ maps Q $\times$ $\Sigma$ (or Q $\times$ ($\Sigma$ $\cup$ \{$\varepsilon$\})) to 2\textsuperscript{Q}.}
\ans{19}{3}{The number of classes equals the number of states of the minimal DFA.}
\ans{20}{1}{DFA minimization merges indistinguishable states.}
\ans{21}{2}{The initial partition separates accepting and non-accepting states.}
\ans{22}{2}{It is the power set of the countable set $\Sigma$*, which is uncountable by Cantor's diagonal argument.}
\ans{23}{2}{Both are true, but A follows from counting: there are countably many TMs and uncountably many languages.}
\ans{24}{1}{A number differing from the ith listed number in the ith digit cannot be in the list.}
\ans{25}{3}{Such a set contains itself if and only if it does not; it showed that naive set theory is inconsistent.}
\ans{26}{4}{Strings can be listed by length; the power set of N is uncountable.}
\ans{27}{3}{1 + 2 + 4 + ... + 2\textsuperscript{n} = 2\textsuperscript{n+1} $-$ 1.}
\ans{28}{1}{The counts follow the Fibonacci sequence: 2, 3, 5, 8, 13 for lengths 1 to 5.}
\ans{29}{2}{A shortest accepted string visits no state twice.}
\ans{30}{3}{Finiteness of a regular language is decidable, using exactly the test in Statement I.}
\ans{31}{2}{The DFA counts length modulo lcm(2, 3) = 6.}
\ans{32}{3}{It is often omitted in drawings to give a partial DFA.}
\ans{33}{3}{Token patterns are regular expressions, which are converted to DFAs (as by Lex).}
\ans{34}{3}{Kleene's theorem: finite automata and regular expressions describe the same languages.}
\ans{35}{1}{This language needs a stack (a PDA).}
\ans{36}{4}{An NFA accepts if at least one computation path ends in an accepting state.}
\ans{37}{4}{States for "last symbol not 1", "last symbol 1" and a dead state after 11.}
\ans{38}{2}{Swapping accepting states complements a DFA (reversing transitions gives the reversal); for an NFA the swap does not work in general.}
\end{answers}

\section{Regular Expressions, Regular Grammars, Closure Properties and the Pumping Lemma}

\begin{question}{H}
Which of the following regular expressions is equivalent to (a + b)*?
\optsshort{(ab)*}{(a*b*)*}{a* + b*}{a*b*}
\end{question}

\begin{question}{H}
Which of the following identities about regular expressions is FALSE?
\optsshort{(a + b)* = a* + b*}{(ab)*a = a(ba)*}{(a*)* = a*}{a*a* = a*}
\end{question}

\begin{question}{H}
By Arden's theorem, the equation R = Q + RP, where P does not contain $\varepsilon$, has the unique solution
\optsshort{R = Q*P}{R = PQ*}{R = (P + Q)*}{R = QP*}
\end{question}

\begin{question}{H}
The regular expression 0*(10*)* denotes
\optslong{strings ending in 0}{strings with an even number of 1s}{all strings over \{0, 1\}}{strings with no two consecutive 1s}
\end{question}

\begin{question}{H}
A regular expression for strings over \{a, b\} that end with "ab" is
\optsshort{ab(a + b)*}{a*b*ab}{(a + b)*ab(a + b)*}{(a + b)*ab}
\end{question}

\begin{question}{M}
Which regular expression describes strings over \{0, 1\} with at least one 1?
\optsshort{(0 + 1)*0}{0*1*}{(01)*}{0*1(0 + 1)*}
\end{question}

\begin{question}{M}
The regular expression (0 + 1)*1(0 + 1)(0 + 1) describes strings in which
\optslong{the third symbol from the left is 1}{the string ends in 111}{the third symbol from the right is 1}{there are three 1s}
\end{question}

\begin{question}{H}
Match List I with List II.\\
List I (Language): A. \{a\textsuperscript{n} b\textsuperscript{n} : n $\geq$ 0\}  B. \{a\textsuperscript{m} b\textsuperscript{n} : m, n $\geq$ 0\}  C. \{a\textsuperscript{n} b\textsuperscript{n} c\textsuperscript{n} : n $\geq$ 0\}  D. \{w w\textsuperscript{R} : w $\in$ \{a, b\}*\}\\
List II (Smallest class containing it): I. Regular  II. Context-sensitive but not context-free  III. Context-free but not regular (deterministic)  IV. Context-free but not deterministic context-free
\optsshort{A-III, B-I, C-II, D-IV}{A-III, B-II, C-I, D-IV}{A-I, B-III, C-II, D-IV}{A-III, B-I, C-IV, D-II}
\end{question}

\begin{question}{H}
Which of the following languages is regular?
\optslong{\{w : w has equal numbers of a's and b's\}}{\{a\textsuperscript{p} : p is prime\}}{\{a\textsuperscript{n} b\textsuperscript{n} : n $\geq$ 1\}}{\{a\textsuperscript{n} b\textsuperscript{n} : 1 $\leq$ n $\leq$ 100\}}
\end{question}

\begin{question}{H}
Which of the following languages is regular?
\optslong{\{a\textsuperscript{n} b\textsuperscript{m} : n > m\}}{\{w $\in$ \{a, b\}* : \#a(w) mod 3 = \#b(w) mod 3\}}{\{w $\in$ \{a, b\}* : \#a(w) = \#b(w)\}}{\{a\textsuperscript{n\textsuperscript{2}} : n $\geq$ 1\}}
\end{question}

\begin{question}{H}
Given below are two statements.\\
Statement I: If a language satisfies the conditions of the pumping lemma, it must be regular.\\
Statement II: The pumping lemma can be used to prove that a language is not regular.
\optslong{Both Statement I and Statement II are true}{Both Statement I and Statement II are false}{Statement I is true but Statement II is false}{Statement I is false but Statement II is true}
\end{question}

\begin{question}{H}
In the pumping lemma for regular languages, a string w with |w| $\geq$ p is split as w = xyz with
\optslong{|y| $\geq$ p}{|y| $\geq$ 1 and |xy| $\leq$ p, and xy\textsuperscript{i} z $\in$ L for all i $\geq$ 0}{y = $\varepsilon$ allowed and |xy| $\leq$ p}{x = $\varepsilon$ always}
\end{question}

\begin{question}{H}
Match List I with List II.\\
List I (Class of languages): A. Context-free languages  B. Deterministic context-free languages  C. Recursively enumerable languages  D. Regular languages\\
List II (An operation under which the class is NOT closed): I. Union  II. Intersection  III. Complement  IV. Infinite union
\optsshort{A-II, B-I, C-III, D-IV}{A-II, B-III, C-I, D-IV}{A-III, B-I, C-II, D-IV}{A-II, B-I, C-IV, D-III}
\end{question}

\begin{question}{H}
Which of the following operations are regular languages NOT closed under?
\optsshort{Reversal}{Union}{Complement}{Infinite union}
\end{question}

\begin{question}{M}
Given below are two statements: one is labelled as Assertion A and the other is labelled as Reason R.\\
Assertion A: The intersection of two regular languages is regular.\\
Reason R: The intersection is accepted by a DFA whose states are the union of the two DFAs' state sets.
\optslong{Both A and R are true and R is the correct explanation of A}{Both A and R are true but R is NOT the correct explanation of A}{A is true but R is false}{A is false but R is true}
\end{question}

\begin{question}{M}
If L1 is regular and L2 is any language, then L1 $\cup$ L2 is
\optslong{regular or non-regular, depending on L2}{always non-regular}{always regular}{always context-free}
\end{question}

\begin{question}{M}
A right-linear grammar has productions of the form
\optslong{any production}{A $\rightarrow$ Bx}{A $\rightarrow$ BC or A $\rightarrow$ a}{A $\rightarrow$ xB or A $\rightarrow$ x, where x is a string of terminals}
\end{question}

\begin{question}{M}
Given below are two statements.\\
Statement I: A grammar mixing left-linear and right-linear productions always generates a regular language.\\
Statement II: Every regular language can be generated by a right-linear grammar.
\optslong{Both Statement I and Statement II are true}{Both Statement I and Statement II are false}{Statement I is true but Statement II is false}{Statement I is false but Statement II is true}
\end{question}

\begin{question}{M}
The regular expression for the language over \{a, b\} in which every string has an even number of a's is
\optsshort{(aa + b)*}{a*b*a*}{(b*ab*ab*)*}{b*(ab*ab*)*}
\end{question}

\begin{question}{M}
Which of the following is the language of the regular expression (a + $\varepsilon$)(b + $\varepsilon$)?
\optsshort{\{$\varepsilon$, ab\}}{\{$\varepsilon$, a, b, ab\}}{\{ab\}}{\{a, b, ab\}}
\end{question}

\begin{question}{M}
The value of $\emptyset$* is
\optsshort{undefined}{\{$\varepsilon$\}}{$\Sigma$*}{$\emptyset$}
\end{question}

\begin{question}{M}
Which of the following problems is decidable for regular languages?\\
A. Membership  B. Emptiness  C. Finiteness  D. Equivalence of two regular languages\\
Choose the correct answer from the options given below:
\optsshort{B and D only}{A, B, C and D}{A and B only}{A, B and C only}
\end{question}

\begin{question}{M}
The reversal of a regular language is
\optslong{context-free but not regular}{never regular}{regular only if the language is finite}{always regular}
\end{question}

\begin{question}{M}
The quotient L1/L2 = \{x : xy $\in$ L1 for some y $\in$ L2\} of a regular language L1 by any language L2 is
\optslong{context-free}{never regular}{regular only if L2 is regular}{always regular}
\end{question}

\begin{question}{M}
Given below are two statements: one is labelled as Assertion A and the other is labelled as Reason R.\\
Assertion A: The language \{a\textsuperscript{p} : p is prime\} is not regular.\\
Reason R: Its complement is finite.
\optslong{Both A and R are true and R is the correct explanation of A}{Both A and R are true but R is NOT the correct explanation of A}{A is true but R is false}{A is false but R is true}
\end{question}

\begin{question}{M}
Which of the following regular expressions denotes the strings over \{0, 1\} in which every 0 is immediately followed by 1?
\optsshort{(1 + 01)*}{(01)*}{1*01*}{(0 + 1)*01}
\end{question}

\begin{question}{M}
The minimum length of the strings generated by the regular expression (a + b)(a + b)(aa + bb)* is
\optsshort{1}{4}{0}{2}
\end{question}

\begin{question}{M}
The regular expression a(a + b)*b + b(a + b)*a describes strings over \{a, b\} that
\optslong{start and end with the same symbol}{start and end with different symbols}{contain ab}{have even length}
\end{question}

\begin{question}{M}
Given below are two statements.\\
Statement I: Every finite language is regular.\\
Statement II: Every regular language is finite.
\optslong{Both Statement I and Statement II are true}{Both Statement I and Statement II are false}{Statement I is true but Statement II is false}{Statement I is false but Statement II is true}
\end{question}

\begin{question}{L}
Regular expressions were introduced by
\optsshort{Stephen Kleene}{John Backus}{Noam Chomsky}{Alan Turing}
\end{question}

\begin{question}{M}
Which of the following languages over \{a, b\} is NOT regular?
\optslong{Strings ending in b}{Strings of length divisible by 4}{Strings with an even number of a's}{Strings in which the number of a's equals the number of b's}
\end{question}

\begin{question}{M}
If L is regular, then which of the following is also regular?\\
A. The set of prefixes of strings in L  B. The set of suffixes of strings in L  C. The complement of L  D. L $\cap$ \{a\textsuperscript{n} b\textsuperscript{n}\}\\
Choose the correct answer from the options given below:
\optsshort{A and D only}{A, B and C only}{B, C and D only}{A, B, C and D}
\end{question}

\begin{question}{M}
Given below are two statements.\\
Statement I: The language \{a\textsuperscript{n} b\textsuperscript{n} : n $\geq$ 0\} $\cup$ \{a\textsuperscript{m} b\textsuperscript{n} : m $\neq$ n\} is regular.\\
Statement II: The union of two non-regular languages is always non-regular.
\optslong{Both Statement I and Statement II are true}{Both Statement I and Statement II are false}{Statement I is true but Statement II is false}{Statement I is false but Statement II is true}
\end{question}

\begin{answers}
\ans{1}{2}{(a*b*)* can generate any string of a's and b's; a* + b* and a*b* cannot generate "ba".}
\ans{2}{1}{a* + b* does not contain mixed strings such as ab.}
\ans{3}{4}{It is used to derive a regular expression from a finite automaton.}
\ans{4}{3}{Any string can be written as a block of 0s followed by blocks each starting with a 1.}
\ans{5}{4}{ab(a + b)* gives strings starting with ab, and (a + b)*ab(a + b)* gives those containing ab.}
\ans{6}{4}{The first 1 can be preceded by any number of 0s and followed by anything.}
\ans{7}{3}{Exactly two symbols follow the chosen 1.}
\ans{8}{1}{\{w w\textsuperscript{R}\} needs non-determinism to guess the middle; \{w c w\textsuperscript{R}\} is deterministic.}
\ans{9}{4}{Any finite language is regular.}
\ans{10}{2}{A DFA tracks (\#a $-$ \#b) mod 3 with 3 states.}
\ans{11}{4}{Some non-regular languages satisfy the pumping conditions, so the lemma is only a necessary condition.}
\ans{12}{2}{p is the pumping length (at most the number of DFA states).}
\ans{13}{1}{RE languages are closed under intersection and union, DCFLs under complement, and regular languages under complement, which rules out each distractor.}
\ans{14}{4}{An infinite union of the finite languages \{a\textsuperscript{n} b\textsuperscript{n}\} gives the non-regular \{a\textsuperscript{n} b\textsuperscript{n} : n $\geq$ 0\}.}
\ans{15}{3}{The product construction uses pairs of states, running both DFAs in parallel.}
\ans{16}{1}{For example, with L1 = $\emptyset$ the result is L2, which may be non-regular.}
\ans{17}{4}{Right-linear and left-linear grammars both generate exactly the regular languages.}
\ans{18}{4}{Mixing them can generate \{a\textsuperscript{n} b\textsuperscript{n}\}: S $\rightarrow$ aA, A $\rightarrow$ Sb, S $\rightarrow$ $\varepsilon$.}
\ans{19}{4}{After any leading b's, a's come in pairs separated by any number of b's.}
\ans{20}{2}{Each factor contributes its symbol or nothing.}
\ans{21}{2}{The Kleene star of any language contains $\varepsilon$; $\varepsilon$* is also \{$\varepsilon$\}.}
\ans{22}{2}{All of these are decidable using DFAs.}
\ans{23}{4}{Reverse the edges of an NFA and swap the start and accepting states.}
\ans{24}{4}{It is a subset of the prefixes of L1 and can be accepted by the DFA of L1 with new accepting states.}
\ans{25}{3}{A is true (by the pumping lemma), but its complement \{a\textsuperscript{n} : n not prime\} is infinite.}
\ans{26}{1}{Each occurrence of 0 is part of a 01 block.}
\ans{27}{4}{The two leading symbols are mandatory and the starred part can be empty.}
\ans{28}{2}{Each alternative starts with one symbol and ends with the other.}
\ans{29}{3}{a* is regular and infinite.}
\ans{30}{1}{Kleene's theorem shows they describe exactly the languages of finite automata.}
\ans{31}{4}{Equality of counts requires unbounded counting.}
\ans{32}{2}{Intersecting with the non-regular \{a\textsuperscript{n} b\textsuperscript{n}\} need not give a regular language.}
\ans{33}{3}{The union in Statement I is a*b*, which is regular, so Statement II is false.}
\end{answers}

\section{Context-Free Grammars, Derivations, Ambiguity, Chomsky and Greibach Normal Forms}

\begin{question}{H}
The grammar E $\rightarrow$ E + E | E * E | id generates the string id + id * id with how many distinct parse trees?
\optsshort{1}{4}{3}{2}
\end{question}

\begin{question}{H}
With the grammar E $\rightarrow$ E + E | E * E | id, the number of parse trees for id + id * id + id (three operators) is
\optsshort{6}{5}{3}{4}
\end{question}

\begin{question}{M}
With the grammar S $\rightarrow$ SS | a, the number of parse trees for the string aaaa is
\optsshort{14}{3}{4}{5}
\end{question}

\begin{question}{H}
A grammar is ambiguous if
\optslong{it has left recursion}{it has $\varepsilon$-productions}{it generates an infinite language}{some string has more than one leftmost derivation (or more than one parse tree)}
\end{question}

\begin{question}{H}
Given below are two statements.\\
Statement I: Whether an arbitrary context-free grammar is ambiguous is undecidable.\\
Statement II: Every context-free language has an unambiguous grammar.
\optslong{Both Statement I and Statement II are true}{Both Statement I and Statement II are false}{Statement I is true but Statement II is false}{Statement I is false but Statement II is true}
\end{question}

\begin{question}{H}
Which of the following languages is inherently ambiguous?
\optslong{\{a\textsuperscript{i} b\textsuperscript{j} c\textsuperscript{k} : i = j or j = k\}}{\{a\textsuperscript{n} b\textsuperscript{m} : n $\neq$ m\}}{\{w w\textsuperscript{R} : w $\in$ \{a, b\}*\}}{\{a\textsuperscript{n} b\textsuperscript{n} : n $\geq$ 0\}}
\end{question}

\begin{question}{H}
In Chomsky normal form, every production is of the form
\optslong{A $\rightarrow$ $\alpha$ with |$\alpha$| $\leq$ 2}{A $\rightarrow$ BC or A $\rightarrow$ a (plus S $\rightarrow$ $\varepsilon$ if $\varepsilon$ is in the language)}{A $\rightarrow$ a$\alpha$ with $\alpha$ a string of variables}{A $\rightarrow$ aB or A $\rightarrow$ a}
\end{question}

\begin{question}{H}
If a grammar is in Chomsky normal form, the number of steps in a derivation of a string of length n $\geq$ 1 is
\optsshort{2n}{n}{2n $-$ 1}{n + 1}
\end{question}

\begin{question}{H}
In Greibach normal form, every production is of the form
\optslong{A $\rightarrow$ BC}{A $\rightarrow$ a only}{A $\rightarrow$ a$\alpha$, where a is a terminal and $\alpha$ is a (possibly empty) string of variables}{A $\rightarrow$ $\alpha$a}
\end{question}

\begin{question}{H}
Arrange the steps of converting a context-free grammar to Chomsky normal form in a correct order.\\
A. Eliminate unit productions  B. Eliminate $\varepsilon$-productions  C. Break right sides longer than two into binary productions and replace terminals in long right sides by variables  D. Eliminate useless symbols
\optsshort{A, B, C, D}{D, C, B, A}{C, B, A, D}{B, A, D, C}
\end{question}

\begin{question}{M}
Given below are two statements: one is labelled as Assertion A and the other is labelled as Reason R.\\
Assertion A: In a CNF grammar, the parse tree of a string of length n has height at least about log2 n.\\
Reason R: Every internal node of a CNF parse tree that is not directly above a terminal has exactly two children.
\optslong{Both A and R are true and R is the correct explanation of A}{Both A and R are true but R is NOT the correct explanation of A}{A is true but R is false}{A is false but R is true}
\end{question}

\begin{question}{H}
The CYK algorithm decides membership of a string of length n for a grammar in Chomsky normal form in time
\optsshort{O(n\textsuperscript{3})}{O(n)}{O(2\textsuperscript{n})}{O(n\textsuperscript{2})}
\end{question}

\begin{question}{H}
Left recursion A $\rightarrow$ A$\alpha$ | $\beta$ is removed by replacing it with
\optsshort{A $\rightarrow$ A$\beta$ | $\alpha$}{A $\rightarrow$ $\beta\alpha$*}{A $\rightarrow$ $\beta$A', A' $\rightarrow$ $\alpha$A' | $\varepsilon$}{A $\rightarrow$ $\alpha$A', A' $\rightarrow$ $\beta$A' | $\varepsilon$}
\end{question}

\begin{question}{M}
Left factoring is used to transform a grammar so that
\optslong{left recursion is removed}{a predictive parser can choose a production with one symbol of lookahead when alternatives share a common prefix}{the grammar becomes ambiguous}{$\varepsilon$-productions are removed}
\end{question}

\begin{question}{M}
Given below are two statements.\\
Statement I: Removing useless symbols changes the language generated by the grammar.\\
Statement II: A variable is useless if it is not reachable from the start symbol or derives no terminal string.
\optslong{Both Statement I and Statement II are true}{Both Statement I and Statement II are false}{Statement I is true but Statement II is false}{Statement I is false but Statement II is true}
\end{question}

\begin{question}{M}
A unit production is a production of the form
\optslong{A $\rightarrow$ BC}{A $\rightarrow$ $\varepsilon$}{A $\rightarrow$ B, where A and B are variables}{A $\rightarrow$ a}
\end{question}

\begin{question}{M}
The grammar S $\rightarrow$ aSb | $\varepsilon$ generates
\optsshort{\{a\textsuperscript{n} b\textsuperscript{m} : n, m $\geq$ 0\}}{\{(ab)\textsuperscript{n} : n $\geq$ 0\}}{\{a\textsuperscript{n} b\textsuperscript{n} : n $\geq$ 0\}}{\{a\textsuperscript{n} b\textsuperscript{n} : n $\geq$ 1\}}
\end{question}

\begin{question}{M}
The grammar S $\rightarrow$ aSa | bSb | a | b | $\varepsilon$ generates
\optslong{strings with equal numbers of a's and b's}{strings of even length only}{all palindromes over \{a, b\}}{all strings over \{a, b\}}
\end{question}

\begin{question}{M}
The grammar S $\rightarrow$ aSbS | bSaS | $\varepsilon$ generates
\optslong{palindromes}{strings with more a's than b's}{strings with equal numbers of a's and b's}{a\textsuperscript{n} b\textsuperscript{n} only}
\end{question}

\begin{question}{M}
The dangling-else ambiguity in "if E then if E then S else S" is usually resolved by
\optslong{reporting an error}{ignoring the else}{associating each else with the outermost if}{associating each else with the nearest unmatched then (if)}
\end{question}

\begin{question}{M}
To make the expression grammar E $\rightarrow$ E + E | E * E | id unambiguous, giving * higher precedence than + and making both left-associative, we write
\optslong{E $\rightarrow$ E * T | T, T $\rightarrow$ T + F | F, F $\rightarrow$ id}{E $\rightarrow$ T + E | T, T $\rightarrow$ F * T | F, F $\rightarrow$ id}{E $\rightarrow$ E + T | T, T $\rightarrow$ T * F | F, F $\rightarrow$ id}{E $\rightarrow$ id + E | id * E | id}
\end{question}

\begin{question}{M}
Given below are two statements: one is labelled as Assertion A and the other is labelled as Reason R.\\
Assertion A: The grammar E $\rightarrow$ E $-$ E | id generates the right language for subtraction but is unsuitable for compilers.\\
Reason R: The grammar is ambiguous, so id $-$ id $-$ id could be grouped as (id $-$ id) $-$ id or id $-$ (id $-$ id), which give different values.
\optslong{Both A and R are true and R is the correct explanation of A}{Both A and R are true but R is NOT the correct explanation of A}{A is true but R is false}{A is false but R is true}
\end{question}

\begin{question}{M}
A leftmost derivation is one in which
\optslong{terminals are replaced first}{the start symbol is never used}{the leftmost variable is replaced at each step}{the rightmost variable is replaced at each step}
\end{question}

\begin{question}{M}
Given below are two statements.\\
Statement I: Each parse tree corresponds to exactly one leftmost derivation.\\
Statement II: A string with two different rightmost derivations always has two different parse trees.
\optslong{Both Statement I and Statement II are true}{Both Statement I and Statement II are false}{Statement I is true but Statement II is false}{Statement I is false but Statement II is true}
\end{question}

\begin{question}{M}
The yield of a parse tree is
\optslong{the height of the tree}{the number of nodes}{the string of leaf labels read from left to right}{the root label}
\end{question}

\begin{question}{M}
A context-free grammar G = (V, T, P, S) is in Chomsky normal form with 3 variables. The maximum number of distinct productions of the form A $\rightarrow$ BC is
\optsshort{9}{81}{27}{18}
\end{question}

\begin{question}{M}
Which of the following grammars is left-recursive?
\optsshort{S $\rightarrow$ Sa | b}{S $\rightarrow$ ab}{S $\rightarrow$ aSb | $\varepsilon$}{S $\rightarrow$ aS | b}
\end{question}

\begin{question}{M}
Given below are two statements.\\
Statement I: Every context-free grammar is a regular grammar.\\
Statement II: Every context-free language is regular.
\optslong{Both Statement I and Statement II are true}{Both Statement I and Statement II are false}{Statement I is true but Statement II is false}{Statement I is false but Statement II is true}
\end{question}

\begin{question}{L}
BNF (Backus-Naur Form) is a notation used to describe
\optslong{context-free grammars of programming languages}{regular expressions}{machine code}{lexical tokens only}
\end{question}

\begin{question}{M}
Given below are two statements: one is labelled as Assertion A and the other is labelled as Reason R.\\
Assertion A: Converting a grammar to Chomsky normal form always removes left recursion.\\
Reason R: In Greibach normal form every production begins with a terminal.
\optslong{Both A and R are true and R is the correct explanation of A}{Both A and R are true but R is NOT the correct explanation of A}{A is true but R is false}{A is false but R is true}
\end{question}

\begin{question}{M}
The grammar S $\rightarrow$ aB | bA, A $\rightarrow$ a | aS | bAA, B $\rightarrow$ b | bS | aBB generates
\optslong{a\textsuperscript{n} b\textsuperscript{n}}{palindromes}{strings with equal numbers of a's and b's (non-empty)}{strings ending with a}
\end{question}

\begin{question}{M}
Which of the following is the correct leftmost derivation of id + id using E $\rightarrow$ E + T | T, T $\rightarrow$ id?
\optslong{E $\Rightarrow$ T $\Rightarrow$ id}{E $\Rightarrow$ E + T $\Rightarrow$ T + T $\Rightarrow$ id + T $\Rightarrow$ id + id}{E $\Rightarrow$ E + T $\Rightarrow$ E + id $\Rightarrow$ T + id $\Rightarrow$ id + id}{E $\Rightarrow$ E + T $\Rightarrow$ id + id}
\end{question}

\begin{answers}
\ans{1}{4}{One tree groups (id + id) * id and the other id + (id * id), so the grammar is ambiguous.}
\ans{2}{2}{The number of ways to bracket k binary operators is the Catalan number C(k) = C(3) = 5.}
\ans{3}{4}{It is the Catalan number C(3) = 5.}
\ans{4}{4}{Equivalently, some string has more than one rightmost derivation.}
\ans{5}{3}{Some CFLs, such as \{a\textsuperscript{i} b\textsuperscript{j} c\textsuperscript{k} : i = j or j = k\}, are inherently ambiguous.}
\ans{6}{1}{Strings with i = j = k have two essentially different derivations in every grammar.}
\ans{7}{2}{B and C are variables and a is a terminal.}
\ans{8}{3}{n $-$ 1 applications of A $\rightarrow$ BC and n applications of A $\rightarrow$ a.}
\ans{9}{3}{Each step of a derivation produces exactly one terminal, so a string of length n needs n steps.}
\ans{10}{4}{Removing $\varepsilon$-productions can create unit productions, and both can create useless symbols, so this order is safe.}
\ans{11}{1}{The tree is essentially binary, which is used in the pumping lemma for CFLs.}
\ans{12}{1}{It fills a triangular table over all substrings and split points.}
\ans{13}{3}{Both generate $\beta$ $\alpha$*, but the new grammar is right-recursive and suits top-down parsing.}
\ans{14}{2}{A $\rightarrow$ $\alpha\beta$1 | $\alpha\beta$2 becomes A $\rightarrow$ $\alpha$A', A' $\rightarrow$ $\beta$1 | $\beta$2.}
\ans{15}{4}{Useless symbols take part in no derivation of a terminal string, so removing them does not change the language.}
\ans{16}{3}{Unit productions are removed by substituting the productions of B.}
\ans{17}{3}{Each application of S $\rightarrow$ aSb adds one a and one b.}
\ans{18}{3}{$\varepsilon$ gives even-length palindromes and a or b gives odd-length ones.}
\ans{19}{3}{Every production adds one a and one b, and the grammar can generate any balanced arrangement.}
\ans{20}{4}{Parser generators such as Yacc resolve the shift-reduce conflict in favour of shifting.}
\ans{21}{3}{Precedence is encoded by levels; left recursion gives left associativity.}
\ans{22}{1}{Compilers need a grammar that encodes associativity.}
\ans{23}{3}{Top-down parsers construct leftmost derivations; bottom-up parsers construct rightmost derivations in reverse.}
\ans{24}{1}{Parse trees correspond one-to-one with leftmost derivations, and also with rightmost derivations.}
\ans{25}{3}{For a complete parse tree it is a terminal string in the language.}
\ans{26}{3}{3 choices for each of A, B and C: 3\textsuperscript{3} = 27.}
\ans{27}{1}{A production whose right side begins with its own left-hand variable is left-recursive.}
\ans{28}{2}{Regular grammars are a restricted kind of CFG, and \{a\textsuperscript{n} b\textsuperscript{n}\} is context-free but not regular.}
\ans{29}{1}{It was first used to define ALGOL 60.}
\ans{30}{4}{CNF productions such as A $\rightarrow$ AB can still be left-recursive; it is GNF that removes left recursion, because of R.}
\ans{31}{3}{A derives strings with one more a than b, and B strings with one more b than a.}
\ans{32}{2}{The second option is a rightmost derivation.}
\end{answers}

\section{Pushdown Automata and Properties of Context-Free Languages}

\begin{question}{H}
Match List I with List II.\\
List I (Language class): A. Regular  B. Context-free  C. Context-sensitive  D. Recursively enumerable\\
List II (Accepting automaton): I. Linear-bounded automaton  II. Turing machine  III. Finite automaton  IV. Non-deterministic pushdown automaton
\optsshort{A-III, B-IV, C-II, D-I}{A-III, B-I, C-IV, D-II}{A-III, B-IV, C-I, D-II}{A-IV, B-III, C-I, D-II}
\end{question}

\begin{question}{H}
Given below are two statements.\\
Statement I: Deterministic PDAs accept exactly the same class of languages as non-deterministic PDAs.\\
Statement II: The language \{w w\textsuperscript{R} : w $\in$ \{a, b\}*\} can be accepted by a non-deterministic PDA but not by any deterministic PDA.
\optslong{Both Statement I and Statement II are true}{Both Statement I and Statement II are false}{Statement I is true but Statement II is false}{Statement I is false but Statement II is true}
\end{question}

\begin{question}{H}
Which of the following languages is a deterministic context-free language?
\optsshort{\{w w : w $\in$ \{a, b\}*\}}{\{w w\textsuperscript{R} : w $\in$ \{a, b\}*\}}{\{a\textsuperscript{n} b\textsuperscript{n} c\textsuperscript{n} : n $\geq$ 0\}}{\{w c w\textsuperscript{R} : w $\in$ \{a, b\}*\}}
\end{question}

\begin{question}{H}
Which of the following languages is NOT context-free?
\optsshort{\{a\textsuperscript{n} b\textsuperscript{n} c\textsuperscript{n} : n $\geq$ 0\}}{\{a\textsuperscript{n} b\textsuperscript{m} c\textsuperscript{n} : n, m $\geq$ 0\}}{\{a\textsuperscript{n} b\textsuperscript{n} : n $\geq$ 0\}}{\{a\textsuperscript{n} b\textsuperscript{n} c\textsuperscript{m} d\textsuperscript{m} : n, m $\geq$ 0\}}
\end{question}

\begin{question}{H}
Which of the following languages is NOT context-free?
\optsshort{\{a\textsuperscript{n} b\textsuperscript{n} c\textsuperscript{m} d\textsuperscript{m}\}}{\{a\textsuperscript{n} b\textsuperscript{m} c\textsuperscript{n} d\textsuperscript{m}\}}{\{a\textsuperscript{n} b\textsuperscript{m} c\textsuperscript{m} d\textsuperscript{n}\}}{\{a\textsuperscript{m} b\textsuperscript{n} : m $\neq$ n\}}
\end{question}

\begin{question}{H}
Which of the following languages is NOT context-free?
\optslong{\{w : w has equal numbers of a's and b's\}}{\{w w : w $\in$ \{a, b\}*\}}{\{w w\textsuperscript{R} : w $\in$ \{a, b\}*\}}{\{a\textsuperscript{i} b\textsuperscript{j} : i $\leq$ j\}}
\end{question}

\begin{question}{H}
Context-free languages are closed under which of the following operations?\\
A. Union  B. Intersection  C. Complement  D. Kleene star  E. Intersection with a regular language\\
Choose the correct answer from the options given below:
\optsshort{A, B, C, D and E}{B, C and E only}{A, D and E only}{A, B and D only}
\end{question}

\begin{question}{H}
Given below are two statements: one is labelled as Assertion A and the other is labelled as Reason R.\\
Assertion A: Context-free languages are not closed under intersection.\\
Reason R: Context-free languages are closed under union.
\optslong{Both A and R are true and R is the correct explanation of A}{Both A and R are true but R is NOT the correct explanation of A}{A is true but R is false}{A is false but R is true}
\end{question}

\begin{question}{M}
Given below are two statements.\\
Statement I: Deterministic context-free languages are closed under union.\\
Statement II: Deterministic context-free languages are closed under complement.
\optslong{Both Statement I and Statement II are true}{Both Statement I and Statement II are false}{Statement I is true but Statement II is false}{Statement I is false but Statement II is true}
\end{question}

\begin{question}{H}
A PDA can accept a language by
\optslong{final state only}{either final state or empty stack; both give the same class of languages for non-deterministic PDAs}{neither}{empty stack only}
\end{question}

\begin{question}{H}
A PDA with two stacks is equivalent in power to a
\optsshort{linear-bounded automaton}{Turing machine}{one-stack PDA}{finite automaton}
\end{question}

\begin{question}{M}
The transition function of a non-deterministic PDA maps
\optslong{Q $\times$ ($\Sigma$ $\cup$ \{$\varepsilon$\}) $\times$ $\Gamma$ to finite subsets of Q $\times$ $\Gamma$*}{Q $\times$ $\Sigma$ $\times$ $\Gamma$ to Q $\times$ $\Gamma$*}{Q $\times$ $\Sigma$ to Q}{Q $\times$ $\Gamma$ to $\Sigma$}
\end{question}

\begin{question}{M}
In the pumping lemma for context-free languages, a long string z in L is written as z = uvwxy with
\optslong{|vx| $\geq$ 1 and |vwx| $\leq$ p, and u v\textsuperscript{i} w x\textsuperscript{i} y $\in$ L for all i $\geq$ 0}{only v pumped}{|uv| $\leq$ p}{v = x = $\varepsilon$ allowed}
\end{question}

\begin{question}{M}
Which problem is decidable for context-free grammars?
\optslong{Whether L(G) = $\Sigma$*}{Whether two CFGs generate the same language}{Whether a CFG is ambiguous}{Whether L(G) is empty}
\end{question}

\begin{question}{H}
Given below are two statements.\\
Statement I: The equivalence of two context-free grammars is undecidable.\\
Statement II: Whether the intersection of two context-free languages is empty is undecidable.
\optslong{Both Statement I and Statement II are true}{Both Statement I and Statement II are false}{Statement I is true but Statement II is false}{Statement I is false but Statement II is true}
\end{question}

\begin{question}{M}
The intersection of a context-free language and a regular language is
\optslong{always context-free}{never context-free}{always regular}{always context-sensitive but not context-free}
\end{question}

\begin{question}{M}
The complement of a context-free language is
\optslong{always regular}{never context-free}{always context-sensitive (recursive)}{always context-free}
\end{question}

\begin{question}{M}
Given below are two statements: one is labelled as Assertion A and the other is labelled as Reason R.\\
Assertion A: The language \{a\textsuperscript{n} b\textsuperscript{n} : n $\geq$ 0\} is context-free but not regular.\\
Reason R: A PDA can push a symbol for each a and pop one for each b, but a finite automaton cannot count without bound.
\optslong{Both A and R are true and R is the correct explanation of A}{Both A and R are true but R is NOT the correct explanation of A}{A is true but R is false}{A is false but R is true}
\end{question}

\begin{question}{M}
Which of the following is accepted by a deterministic PDA?
\optslong{\{a\textsuperscript{n} b\textsuperscript{n} : n $\geq$ 1\} $\cup$ \{a\textsuperscript{n} b\textsuperscript{2n} : n $\geq$ 1\}}{\{a\textsuperscript{n} b\textsuperscript{n} : n $\geq$ 1\}}{\{w w\textsuperscript{R}\}}{\{a\textsuperscript{i} b\textsuperscript{j} c\textsuperscript{k} : i = j or j = k\}}
\end{question}

\begin{question}{M}
Every context-free language without $\varepsilon$ can be generated by a grammar in
\optslong{Greibach normal form only}{Chomsky normal form only}{both Chomsky and Greibach normal forms}{neither}
\end{question}

\begin{question}{M}
Given below are two statements.\\
Statement I: Every regular language is a deterministic context-free language.\\
Statement II: Every deterministic context-free language is regular.
\optslong{Both Statement I and Statement II are true}{Both Statement I and Statement II are false}{Statement I is true but Statement II is false}{Statement I is false but Statement II is true}
\end{question}

\begin{question}{M}
Context-free languages are closed under
\optsshort{reversal}{set difference}{complement}{intersection}
\end{question}

\begin{question}{M}
The language \{a\textsuperscript{m} b\textsuperscript{n} c\textsuperscript{k} : m = n or n = k\} is
\optslong{context-free and inherently ambiguous}{regular}{deterministic context-free}{not context-free}
\end{question}

\begin{question}{M}
Given below are two statements: one is labelled as Assertion A and the other is labelled as Reason R.\\
Assertion A: The language \{a\textsuperscript{n\textsuperscript{2}} : n $\geq$ 0\} is not context-free.\\
Reason R: The language \{a\textsuperscript{n\textsuperscript{2}} : n $\geq$ 0\} can be generated by a regular grammar.
\optslong{Both A and R are true and R is the correct explanation of A}{Both A and R are true but R is NOT the correct explanation of A}{A is true but R is false}{A is false but R is true}
\end{question}

\begin{question}{M}
A PDA for \{a\textsuperscript{n} b\textsuperscript{n} : n $\geq$ 0\} uses the stack to
\optslong{store the state}{count the a's and match each b with one of them}{reverse the input}{store the whole input}
\end{question}

\begin{question}{L}
The stack of a PDA has
\optslong{unbounded size, but access only at the top}{bounded size}{two ends}{random access}
\end{question}

\begin{question}{M}
Which of the following languages is context-free?
\optsshort{\{a\textsuperscript{p} : p prime\}}{\{a\textsuperscript{i} b\textsuperscript{j} c\textsuperscript{k} : i = j\}}{\{a\textsuperscript{n} b\textsuperscript{n} c\textsuperscript{n}\}}{\{w w\}}
\end{question}

\begin{question}{M}
Given below are two statements.\\
Statement I: Whether a context-free language is regular is decidable.\\
Statement II: Whether a context-free language is finite is decidable.
\optslong{Both Statement I and Statement II are true}{Both Statement I and Statement II are false}{Statement I is true but Statement II is false}{Statement I is false but Statement II is true}
\end{question}

\begin{answers}
\ans{1}{3}{The Chomsky hierarchy and the corresponding machines.}
\ans{2}{4}{DPDAs accept the proper subclass of deterministic CFLs.}
\ans{3}{4}{The marker c tells the DPDA when to stop pushing and start matching.}
\ans{4}{1}{A PDA cannot compare three counts; the pumping lemma for CFLs shows it is not context-free.}
\ans{5}{2}{The crossed dependencies (a with c, b with d) cannot be handled by one stack.}
\ans{6}{2}{\{ww\} needs to compare the first half with the second half in the same order, which a stack cannot do.}
\ans{7}{3}{CFLs are closed under union, concatenation, star, homomorphism and intersection with regular languages, but not under intersection or complement.}
\ans{8}{2}{Both are true; A is shown by \{a\textsuperscript{n} b\textsuperscript{n} c\textsuperscript{m}\} $\cap$ \{a\textsuperscript{m} b\textsuperscript{n} c\textsuperscript{n}\} = \{a\textsuperscript{n} b\textsuperscript{n} c\textsuperscript{n}\}.}
\ans{9}{4}{DCFLs are closed under complement but not under union or intersection.}
\ans{10}{2}{For DPDAs, acceptance by empty stack accepts only languages with the prefix property.}
\ans{11}{2}{Two stacks can simulate the tape of a Turing machine, one holding each side of the head.}
\ans{12}{1}{$\Gamma$ is the stack alphabet.}
\ans{13}{1}{Two parts, v and x, are pumped together.}
\ans{14}{4}{Emptiness, finiteness and membership are decidable for CFLs.}
\ans{15}{1}{Both are undecidable; they are proved by reduction from the Post correspondence problem.}
\ans{16}{1}{A PDA and a DFA can run in parallel (product construction).}
\ans{17}{3}{Every CFL is recursive, and recursive languages are closed under complement, but the complement need not be context-free.}
\ans{18}{1}{This is the standard separating example between Type 3 and Type 2.}
\ans{19}{2}{The union in option 1 requires guessing whether to match each a with one b or two.}
\ans{20}{3}{Normal forms do not reduce generative power.}
\ans{21}{3}{\{a\textsuperscript{n} b\textsuperscript{n}\} is a deterministic CFL that is not regular.}
\ans{22}{1}{Reverse every production of the grammar.}
\ans{23}{1}{It is the union of two CFLs, but strings with m = n = k always have two parse trees.}
\ans{24}{3}{Unary CFLs are regular (Parikh), and \{a\textsuperscript{n\textsuperscript{2}}\} is not regular, so it is not context-free either.}
\ans{25}{2}{It pushes on a and pops on b, accepting when the stack is empty at the end.}
\ans{26}{1}{This restricted access is what limits PDAs compared with Turing machines.}
\ans{27}{2}{c\textsuperscript{k} is free, so the language is \{a\textsuperscript{n} b\textsuperscript{n}\} followed by c*.}
\ans{28}{4}{Regularity of a CFL is undecidable; finiteness can be tested by looking for cycles in the variable graph of a CNF grammar.}
\end{answers}

\section{Turing Machines, Chomsky Hierarchy, Decidability and Unsolvable Problems}

\begin{question}{H}
Match List I with List II.\\
List I (Grammar type in the Chomsky hierarchy): A. Type 0  B. Type 1  C. Type 2  D. Type 3\\
List II (Grammar): I. Regular  II. Unrestricted  III. Context-free  IV. Context-sensitive
\optsshort{A-IV, B-II, C-III, D-I}{A-II, B-IV, C-III, D-I}{A-II, B-IV, C-I, D-III}{A-II, B-III, C-IV, D-I}
\end{question}

\begin{question}{H}
A context-sensitive grammar has productions $\alpha$ $\rightarrow$ $\beta$ with
\optslong{|$\alpha$| $\leq$ |$\beta$| (non-contracting), except possibly S $\rightarrow$ $\varepsilon$}{$\beta$ = $\varepsilon$ allowed everywhere}{$\alpha$ a single terminal}{|$\alpha$| = 1}
\end{question}

\begin{question}{H}
Given below are two statements.\\
Statement I: Every recursive language is recursively enumerable.\\
Statement II: Every recursively enumerable language is recursive.
\optslong{Both Statement I and Statement II are true}{Both Statement I and Statement II are false}{Statement I is true but Statement II is false}{Statement I is false but Statement II is true}
\end{question}

\begin{question}{H}
Given below are two statements: one is labelled as Assertion A and the other is labelled as Reason R.\\
Assertion A: If a language L and its complement are both recursively enumerable, then L is recursive.\\
Reason R: Turing machines for L and its complement can be run in parallel, and exactly one of them accepts any given input.
\optslong{Both A and R are true and R is the correct explanation of A}{Both A and R are true but R is NOT the correct explanation of A}{A is true but R is false}{A is false but R is true}
\end{question}

\begin{question}{H}
The halting problem for Turing machines is
\optslong{decidable}{regular}{not recursively enumerable}{recursively enumerable but undecidable}
\end{question}

\begin{question}{H}
The complement of the halting problem is
\optslong{context-free}{recursively enumerable but not recursive}{not recursively enumerable}{recursive}
\end{question}

\begin{question}{H}
Rice's theorem states that
\optslong{every TM halts on some input}{every non-trivial property of the languages recognised by Turing machines is undecidable}{every property of Turing machines is decidable}{the halting problem is decidable for LBAs}
\end{question}

\begin{question}{H}
Which of the following problems is decidable?
\optslong{Whether two TMs accept the same language}{Whether a TM accepts a regular language}{Whether a TM halts on the empty input}{Whether a DFA accepts a given string}
\end{question}

\begin{question}{H}
Match List I with List II.\\
List I (Problem): A. Membership for context-sensitive languages  B. Emptiness for context-free grammars  C. The halting problem  D. The complement of the halting problem\\
List II (Status): I. Not recursively enumerable  II. Decidable, but PSPACE-complete in general  III. Recursively enumerable but not recursive  IV. Decidable in polynomial time
\optsshort{A-IV, B-II, C-III, D-I}{A-II, B-III, C-IV, D-I}{A-II, B-IV, C-I, D-III}{A-II, B-IV, C-III, D-I}
\end{question}

\begin{question}{H}
The Post correspondence problem (PCP) is
\optslong{undecidable in general, but decidable for a one-letter alphabet}{decidable for any alphabet}{in P}{context-free}
\end{question}

\begin{question}{M}
A Post correspondence problem instance has the dominoes (top, bottom): A = (b, bbb), B = (babbb, ba), C = (ba, a). Which sequence is a solution?
\optsshort{A C B}{C A}{A B}{B A A C}
\end{question}

\begin{question}{H}
Given below are two statements.\\
Statement I: A multi-tape Turing machine can recognise languages that no single-tape Turing machine can.\\
Statement II: A non-deterministic Turing machine can be simulated by a deterministic Turing machine.
\optslong{Both Statement I and Statement II are true}{Both Statement I and Statement II are false}{Statement I is true but Statement II is false}{Statement I is false but Statement II is true}
\end{question}

\begin{question}{H}
Which of the following models is NOT equivalent in power to a standard Turing machine?
\optsshort{Multi-tape TM}{Non-deterministic TM}{Linear-bounded automaton}{PDA with two stacks}
\end{question}

\begin{question}{H}
The Church-Turing thesis states that
\optslong{any function that can be computed by an effective procedure can be computed by a Turing machine}{P = NP}{Turing machines are faster than computers}{every problem is decidable}
\end{question}

\begin{question}{M}
A universal Turing machine is one that
\optslong{has only two states}{accepts only regular languages}{takes the encoding of any TM M and an input w and simulates M on w}{accepts every string}
\end{question}

\begin{question}{M}
Given below are two statements: one is labelled as Assertion A and the other is labelled as Reason R.\\
Assertion A: The diagonal language L\_\allowbreak{}d = \{w\_\allowbreak{}i : M\_\allowbreak{}i does not accept w\_\allowbreak{}i\} is recursively enumerable.\\
Reason R: If some TM M\_\allowbreak{}j accepted L\_\allowbreak{}d, then w\_\allowbreak{}j would be in L\_\allowbreak{}d exactly when it is not.
\optslong{Both A and R are true and R is the correct explanation of A}{Both A and R are true but R is NOT the correct explanation of A}{A is true but R is false}{A is false but R is true}
\end{question}

\begin{question}{M}
Which of the following is closed under complement?\\
A. Regular languages  B. Recursive languages  C. Recursively enumerable languages  D. Context-sensitive languages\\
Choose the correct answer from the options given below:
\optsshort{B and C only}{A, B, C and D}{A, B and D only}{A and C only}
\end{question}

\begin{question}{M}
A Turing machine that halts on every input is called a
\optsshort{transducer only}{decider (total TM)}{linear-bounded automaton}{recogniser}
\end{question}

\begin{question}{M}
Arrange the following classes of languages in increasing order of inclusion.\\
A. Context-free  B. Recursive  C. Regular  D. Context-sensitive  E. Recursively enumerable
\optsshort{C, D, A, B, E}{A, C, D, B, E}{C, A, B, D, E}{C, A, D, B, E}
\end{question}

\begin{question}{M}
Given below are two statements.\\
Statement I: Every recursive language is context-sensitive.\\
Statement II: Every context-sensitive language is recursive.
\optslong{Both Statement I and Statement II are true}{Both Statement I and Statement II are false}{Statement I is true but Statement II is false}{Statement I is false but Statement II is true}
\end{question}

\begin{question}{M}
Which of the following problems about a Turing machine M is decidable?
\optslong{Whether M halts on all inputs}{Whether L(M) is regular}{Whether L(M) is empty}{Whether M has more than 10 states}
\end{question}

\begin{question}{M}
Which problem is undecidable for context-free grammars?
\optsshort{Membership}{Whether L(G) = $\Sigma$*}{Finiteness}{Emptiness}
\end{question}

\begin{question}{M}
If problem A reduces to problem B (A $\leq$ B) and A is undecidable, then
\optsshort{nothing can be concluded}{A is decidable}{B is decidable}{B is undecidable}
\end{question}

\begin{question}{M}
A Turing machine that computes the function f(n) = n + 1 on unary input 1\textsuperscript{n} must
\optslong{erase all symbols}{reverse the input}{never halt}{move to the end of the input and write one more 1}
\end{question}

\begin{question}{M}
Given below are two statements.\\
Statement I: Every Turing machine halts on every input.\\
Statement II: The set of strings on which a given TM halts is always recursive.
\optslong{Both Statement I and Statement II are true}{Both Statement I and Statement II are false}{Statement I is true but Statement II is false}{Statement I is false but Statement II is true}
\end{question}

\begin{question}{M}
A linear-bounded automaton is a Turing machine whose tape head
\optslong{has two tapes}{can move only to the right}{has no tape}{stays within the portion of the tape containing the input}
\end{question}

\begin{question}{M}
Given below are two statements: one is labelled as Assertion A and the other is labelled as Reason R.\\
Assertion A: The emptiness problem for linear-bounded automata is undecidable.\\
Reason R: The membership problem for linear-bounded automata is undecidable.
\optslong{Both A and R are true and R is the correct explanation of A}{Both A and R are true but R is NOT the correct explanation of A}{A is true but R is false}{A is false but R is true}
\end{question}

\begin{question}{M}
The number of Turing machines (with a fixed tape alphabet) is
\optsshort{zero}{finite}{countably infinite}{uncountable}
\end{question}

\begin{question}{L}
The Turing machine model was proposed by Alan Turing in
\optsshort{1971}{1956}{1936}{1926}
\end{question}

\begin{question}{M}
In the standard TM, one move consists of
\optslong{reading the whole input at once}{reading a symbol, writing a symbol, moving the head one cell left or right, and changing state}{reading a symbol only}{popping a stack}
\end{question}

\begin{question}{M}
Given below are two statements.\\
Statement I: The class P is closed under complement.\\
Statement II: The class NP is known to be closed under complement.
\optslong{Both Statement I and Statement II are true}{Both Statement I and Statement II are false}{Statement I is true but Statement II is false}{Statement I is false but Statement II is true}
\end{question}

\begin{answers}
\ans{1}{2}{Type 0 $\supset$ Type 1 $\supset$ Type 2 $\supset$ Type 3.}
\ans{2}{1}{Context-sensitive languages are accepted by linear-bounded automata.}
\ans{3}{3}{The halting problem is recursively enumerable but not recursive.}
\ans{4}{1}{The combined machine always halts with the right answer.}
\ans{5}{4}{A universal TM can simulate and accept when the machine halts, but no TM decides it.}
\ans{6}{3}{If both it and HALT were RE, HALT would be decidable.}
\ans{7}{2}{For example, whether L(M) is empty, finite or regular is undecidable.}
\ans{8}{4}{Membership, emptiness and equivalence are decidable for DFAs.}
\ans{9}{4}{If both HALT and its complement were recursively enumerable, HALT would be decidable.}
\ans{10}{1}{PCP is the standard tool for proving undecidability results about CFGs.}
\ans{11}{4}{B A A C gives babbb$\cdot$b$\cdot$b$\cdot$ba = babbbbbba on top and ba$\cdot$bbb$\cdot$bbb$\cdot$a = babbbbbba on the bottom.}
\ans{12}{4}{Multi-tape TMs add speed (at most quadratic simulation overhead) but no power; non-determinism adds no power either.}
\ans{13}{3}{An LBA can use only the tape space occupied by its input and accepts exactly the context-sensitive languages.}
\ans{14}{1}{It is a thesis about the informal notion of computability, not a provable theorem.}
\ans{15}{3}{It is the theoretical model of a stored-program computer.}
\ans{16}{4}{R is the contradiction that proves L\_\allowbreak{}d is not recursively enumerable.}
\ans{17}{3}{CSLs are closed under complement by the Immerman-Szelepcs\'enyi theorem; RE languages are not.}
\ans{18}{2}{The languages of deciders are the recursive languages.}
\ans{19}{4}{Regular $\subset$ CFL $\subset$ CSL $\subset$ recursive $\subset$ RE, all inclusions proper.}
\ans{20}{4}{There are recursive languages that are not context-sensitive (by diagonalization).}
\ans{21}{4}{It is a property of the machine's description, not of its language, so Rice's theorem does not apply.}
\ans{22}{2}{Universality of CFGs is undecidable.}
\ans{23}{4}{A decider for B would give a decider for A.}
\ans{24}{4}{Such simple constructions are typical exam exercises.}
\ans{25}{2}{A TM may loop forever, and its halting set is recursively enumerable but not necessarily recursive.}
\ans{26}{4}{LBAs accept exactly the context-sensitive languages.}
\ans{27}{3}{LBA membership is decidable (finitely many configurations), while LBA emptiness is undecidable.}
\ans{28}{3}{Each TM has a finite description that can be encoded as a string.}
\ans{29}{3}{His paper "On Computable Numbers" also introduced the universal machine.}
\ans{30}{2}{$\delta$: Q $\times$ $\Gamma$ $\rightarrow$ Q $\times$ $\Gamma$ $\times$ \{L, R\}.}
\ans{31}{3}{Whether NP = co-NP is an open problem.}
\end{answers}

\section{Phases of a Compiler and Lexical Analysis}

\begin{question}{H}
Arrange the phases of a compiler in the order in which they process the source program.\\
A. Intermediate code generation  B. Syntax analysis  C. Code optimization  D. Lexical analysis  E. Semantic analysis
\optsshort{D, E, B, A, C}{D, B, A, E, C}{B, D, E, A, C}{D, B, E, A, C}
\end{question}

\begin{question}{H}
Match List I with List II.\\
List I (Compiler phase): A. Lexical analysis  B. Syntax analysis  C. Semantic analysis  D. Code generation\\
List II (Main activity): I. Type checking and building attribute information  II. Grouping characters into tokens  III. Mapping intermediate code to target machine instructions and allocating registers  IV. Building a parse tree from tokens according to the grammar
\optsshort{A-II, B-I, C-IV, D-III}{A-II, B-IV, C-I, D-III}{A-II, B-IV, C-III, D-I}{A-IV, B-II, C-I, D-III}
\end{question}

\begin{question}{H}
The number of tokens in the C statement printf("i = \%d, \&i = \%x", i, \&i); is
\optsshort{9}{13}{8}{10}
\end{question}

\begin{question}{H}
The number of tokens in the C declaration int a[10], *p = \&a[0]; is
\optsshort{14}{13}{16}{15}
\end{question}

\begin{question}{M}
The number of tokens in the C statement x = y++ + ++z; is
\optsshort{8}{7}{6}{9}
\end{question}

\begin{question}{H}
Given below are two statements.\\
Statement I: A lexeme is a sequence of characters in the source program that matches the pattern for a token.\\
Statement II: A token is a pair consisting of a token name and an optional attribute value.
\optslong{Both Statement I and Statement II are true}{Both Statement I and Statement II are false}{Statement I is true but Statement II is false}{Statement I is false but Statement II is true}
\end{question}

\begin{question}{H}
Which of the following is usually detected by the lexical analyzer?
\optslong{An undeclared variable}{Type mismatch in an expression}{A missing semicolon}{An illegal character such as @ in a C program outside a string}
\end{question}

\begin{question}{H}
Match List I with List II.\\
List I (Error): A. The character @ appearing outside a string in a C program  B. Unbalanced parentheses  C. Adding an integer to an array  D. Division by a zero value read from the input\\
List II (Where it is detected): I. Semantic analysis  II. Syntax analysis  III. Lexical analysis  IV. Only at run time
\optsshort{A-III, B-I, C-II, D-IV}{A-III, B-II, C-IV, D-I}{A-II, B-III, C-I, D-IV}{A-III, B-II, C-I, D-IV}
\end{question}

\begin{question}{M}
The symbol table is accessed by
\optsshort{only the parser}{most phases of the compiler}{only the code generator}{only the lexical analyzer}
\end{question}

\begin{question}{H}
Given below are two statements: one is labelled as Assertion A and the other is labelled as Reason R.\\
Assertion A: The front end of a compiler is largely independent of the target machine.\\
Reason R: The front end analyses the source program and produces an intermediate representation, leaving machine-specific work to the back end.
\optslong{Both A and R are true and R is the correct explanation of A}{Both A and R are true but R is NOT the correct explanation of A}{A is true but R is false}{A is false but R is true}
\end{question}

\begin{question}{M}
Lex (or Flex) generates
\optslong{a parser from a context-free grammar}{an optimizer}{an assembler}{a lexical analyzer from regular-expression patterns}
\end{question}

\begin{question}{M}
A Lex specification consists of three sections separated by \%\%. They are
\optslong{tokens, grammar and actions}{header, body and footer}{declarations, translation rules, and auxiliary functions}{states, transitions and outputs}
\end{question}

\begin{question}{M}
When two Lex patterns match the input, Lex chooses
\optslong{the longest match, and among equal-length matches the pattern listed first}{the shorter match}{the pattern listed last}{a random pattern}
\end{question}

\begin{question}{M}
Input buffering with sentinels in a lexical analyzer is used to
\optslong{detect syntax errors}{store tokens}{reduce memory}{reduce the number of tests for the end of the buffer for each character read}
\end{question}

\begin{question}{M}
Given below are two statements.\\
Statement I: A single-pass compiler performs all phases in one scan of the source program.\\
Statement II: A multi-pass compiler always produces faster target code than a single-pass compiler.
\optslong{Both Statement I and Statement II are true}{Both Statement I and Statement II are false}{Statement I is true but Statement II is false}{Statement I is false but Statement II is true}
\end{question}

\begin{question}{M}
The task of recognising keywords such as "if" and "while" is usually done by
\optslong{the parser only}{the optimizer}{treating them as reserved words, entered in the symbol table or matched by separate patterns}{the code generator}
\end{question}

\begin{question}{M}
Panic-mode error recovery in a lexical analyzer
\optslong{deletes successive characters until a well-formed token is found}{rewrites the source file}{asks the user}{stops compilation immediately}
\end{question}

\begin{question}{M}
Given below are two statements: one is labelled as Assertion A and the other is labelled as Reason R.\\
Assertion A: Lexical analysis is separated from syntax analysis.\\
Reason R: Recognising tokens requires the full power of context-free grammars.
\optslong{Both A and R are true and R is the correct explanation of A}{Both A and R are true but R is NOT the correct explanation of A}{A is true but R is false}{A is false but R is true}
\end{question}

\begin{question}{M}
Which of the following is NOT a token class in a typical programming language?
\optsshort{Keyword}{Identifier}{Operator}{Comment}
\end{question}

\begin{question}{M}
Match List I with List II.\\
List I (Tool or concept): A. Yacc  B. Lex  C. Preprocessor  D. Assembler\\
List II (Function): I. Converts assembly language into relocatable machine code  II. Generates an LALR parser  III. Generates a lexical analyzer  IV. Handles macros and file inclusion before compilation
\optsshort{A-III, B-II, C-IV, D-I}{A-II, B-III, C-IV, D-I}{A-II, B-IV, C-III, D-I}{A-II, B-III, C-I, D-IV}
\end{question}

\begin{question}{M}
In the Lex pattern [a-zA-Z\_\allowbreak{}][a-zA-Z0-9\_\allowbreak{}]*, the strings matched are
\optslong{identifiers that start with a letter or underscore}{floating-point numbers}{integers}{keywords only}
\end{question}

\begin{question}{M}
Given below are two statements.\\
Statement I: Code optimization is an optional phase of a compiler.\\
Statement II: Intermediate code makes it easier to retarget a compiler to a new machine.
\optslong{Both Statement I and Statement II are true}{Both Statement I and Statement II are false}{Statement I is true but Statement II is false}{Statement I is false but Statement II is true}
\end{question}

\begin{question}{L}
The output of the lexical analyzer is
\optsshort{machine code}{a stream of tokens}{a parse tree}{three-address code}
\end{question}

\begin{question}{M}
Which phase of the compiler is responsible for checking that a variable is declared before it is used?
\optsshort{Code generation}{Lexical analysis}{Semantic analysis}{Syntax analysis}
\end{question}

\begin{question}{M}
A compiler written in language L that runs on machine M and generates code for machine N is described as
\optsshort{an interpreter}{a cross compiler}{a native compiler}{a decompiler}
\end{question}

\begin{question}{M}
Given below are two statements: one is labelled as Assertion A and the other is labelled as Reason R.\\
Assertion A: Lexical analyzers are usually implemented as DFAs.\\
Reason R: Token patterns cannot be described by regular expressions.
\optslong{Both A and R are true and R is the correct explanation of A}{Both A and R are true but R is NOT the correct explanation of A}{A is true but R is false}{A is false but R is true}
\end{question}

\begin{question}{M}
The phases of a compiler that depend mainly on the source language and not on the target machine form the
\optsshort{loader}{front end}{back end}{middle end only}
\end{question}

\begin{question}{M}
The number of tokens in the statement if (a >= b) max = a; else max = b; is
\optsshort{16}{14}{15}{17}
\end{question}

\begin{answers}
\ans{1}{4}{Lexical, syntax and semantic analysis, then intermediate code generation, optimization and target code generation.}
\ans{2}{2}{The analysis phases form the front end and the synthesis phases the back end.}
\ans{3}{4}{printf ( "i = \%d, \&i = \%x" , i , \& i ) ; --- the string literal is a single token.}
\ans{4}{4}{int a [ 10 ] , * p = \& a [ 0 ] ; --- 15 tokens.}
\ans{5}{1}{x = y ++ + ++ z ; --- the lexer uses the longest-match rule for ++.}
\ans{6}{1}{For example the lexeme "count" gives the token (id, pointer to the symbol-table entry).}
\ans{7}{4}{Missing semicolons are syntax errors; type mismatches and undeclared names are semantic errors.}
\ans{8}{4}{A value that is known only at run time cannot be checked by any compiler phase.}
\ans{9}{2}{Entries are created during analysis and used during semantic checking and code generation.}
\ans{10}{1}{This separation allows m front ends and n back ends to give m $\times$ n compilers.}
\ans{11}{4}{Yacc (or Bison) generates LALR(1) parsers.}
\ans{12}{3}{Translation rules have the form "pattern \{ action \}".}
\ans{13}{1}{So "if" is a keyword only if the keyword pattern is listed before the identifier pattern.}
\ans{14}{4}{A special sentinel character marks the end of each buffer half.}
\ans{15}{3}{Extra passes allow more optimization but do not guarantee faster code.}
\ans{16}{3}{Reserved words cannot be used as identifiers.}
\ans{17}{1}{Other options include inserting, deleting, replacing or transposing characters.}
\ans{18}{3}{Tokens are described by regular expressions, which allows simple and fast DFA-based scanners.}
\ans{19}{4}{Comments and white space are stripped by the lexer and not passed to the parser.}
\ans{20}{2}{Language-processing tools.}
\ans{21}{1}{The first character cannot be a digit.}
\ans{22}{1}{Only analysis and code generation are essential.}
\ans{23}{2}{The parser requests tokens one at a time, typically through a getNextToken function.}
\ans{24}{3}{Declaration-before-use is a context-sensitive property not captured by a CFG.}
\ans{25}{2}{A native compiler generates code for the machine it runs on.}
\ans{26}{3}{Token patterns are regular expressions, which is exactly why they can be compiled to DFAs.}
\ans{27}{2}{Lexical, syntax and semantic analysis and intermediate code generation.}
\ans{28}{3}{if ( a >= b ) max = a ; else max = b ; --- 15 tokens.}
\end{answers}

\section{Syntax Analysis: Top-Down Parsing (FIRST, FOLLOW, recursive descent, LL(1))}

\begin{question}{H}
For the grammar E $\rightarrow$ TE', E' $\rightarrow$ +TE' | $\varepsilon$, T $\rightarrow$ FT', T' $\rightarrow$ *FT' | $\varepsilon$, F $\rightarrow$ (E) | id, the set FIRST(E') is
\optsshort{\{(, id\}}{\{+, $\varepsilon$\}}{\{+, ), \$\}}{\{+\}}
\end{question}

\begin{question}{H}
For the grammar E $\rightarrow$ TE', E' $\rightarrow$ +TE' | $\varepsilon$, T $\rightarrow$ FT', T' $\rightarrow$ *FT' | $\varepsilon$, F $\rightarrow$ (E) | id, the set FOLLOW(T) is
\optsshort{\{+, ), \$\}}{\{), \$\}}{\{+, *\}}{\{*, +, ), \$\}}
\end{question}

\begin{question}{H}
For the grammar E $\rightarrow$ TE', E' $\rightarrow$ +TE' | $\varepsilon$, T $\rightarrow$ FT', T' $\rightarrow$ *FT' | $\varepsilon$, F $\rightarrow$ (E) | id, the set FOLLOW(F) is
\optsshort{\{(, id\}}{\{*, \$\}}{\{*, +, ), \$\}}{\{+, ), \$\}}
\end{question}

\begin{question}{H}
For the grammar E $\rightarrow$ TE', E' $\rightarrow$ +TE' | $\varepsilon$, T $\rightarrow$ FT', T' $\rightarrow$ *FT' | $\varepsilon$, F $\rightarrow$ (E) | id, the set FOLLOW(E) is
\optsshort{\{(, id\}}{\{+, ), \$\}}{\{\$\}}{\{), \$\}}
\end{question}

\begin{question}{H}
In the LL(1) parsing table for E $\rightarrow$ TE', E' $\rightarrow$ +TE' | $\varepsilon$, T $\rightarrow$ FT', T' $\rightarrow$ *FT' | $\varepsilon$, F $\rightarrow$ (E) | id, the production E' $\rightarrow$ $\varepsilon$ is placed in the entries
\optsshort{M[E', id]}{M[E', +]}{M[E', )] and M[E', \$]}{M[E', *]}
\end{question}

\begin{question}{H}
A grammar is LL(1) if, for every pair of productions A $\rightarrow$ $\alpha$ | $\beta$,
\optslong{FIRST($\alpha$) and FIRST($\beta$) are disjoint, and if $\beta$ $\Rightarrow$* $\varepsilon$ then FIRST($\alpha$) and FOLLOW(A) are disjoint}{$\alpha$ and $\beta$ have the same length}{A is the start symbol}{the grammar is ambiguous}
\end{question}

\begin{question}{H}
Given below are two statements.\\
Statement I: A left-recursive grammar cannot be LL(1).\\
Statement II: An ambiguous grammar cannot be LL(1).
\optslong{Both Statement I and Statement II are true}{Both Statement I and Statement II are false}{Statement I is true but Statement II is false}{Statement I is false but Statement II is true}
\end{question}

\begin{question}{H}
The grammar S $\rightarrow$ iEtSS' | a, S' $\rightarrow$ eS | $\varepsilon$, E $\rightarrow$ b (dangling else) is not LL(1) because
\optslong{FIRST(S) is empty}{it has no $\varepsilon$-productions}{the entry M[S', e] contains both S' $\rightarrow$ eS and S' $\rightarrow$ $\varepsilon$, since e is in FOLLOW(S')}{it is left-recursive}
\end{question}

\begin{question}{M}
The grammar S $\rightarrow$ AaAb | BbBa, A $\rightarrow$ $\varepsilon$, B $\rightarrow$ $\varepsilon$ is
\optsshort{neither LL(1) nor LR(1)}{LL(1) but not SLR(1)}{SLR(1) but not LL(1)}{both LL(1) and SLR(1)}
\end{question}

\begin{question}{H}
Which of the following grammars is LL(1)?
\optsshort{S $\rightarrow$ A | B, A $\rightarrow$ a, B $\rightarrow$ a}{S $\rightarrow$ aB | bA, A $\rightarrow$ a, B $\rightarrow$ b}{S $\rightarrow$ Sa | b}{S $\rightarrow$ aS | a}
\end{question}

\begin{question}{M}
The grammar S $\rightarrow$ aS | a is not LL(1). The transformation needed is
\optslong{conversion to CNF}{adding an end marker}{removal of left recursion}{left factoring: S $\rightarrow$ aS', S' $\rightarrow$ S | $\varepsilon$}
\end{question}

\begin{question}{H}
A recursive-descent parser
\optslong{uses an LR table}{cannot handle $\varepsilon$-productions}{is a bottom-up parser}{is a top-down parser with one procedure for each non-terminal}
\end{question}

\begin{question}{H}
Given below are two statements: one is labelled as Assertion A and the other is labelled as Reason R.\\
Assertion A: A recursive-descent parser can loop forever on a left-recursive grammar.\\
Reason R: For A $\rightarrow$ A$\alpha$, the procedure for A calls itself before consuming any input.
\optslong{Both A and R are true and R is the correct explanation of A}{Both A and R are true but R is NOT the correct explanation of A}{A is true but R is false}{A is false but R is true}
\end{question}

\begin{question}{H}
A top-down parser produces
\optslong{a rightmost derivation}{a leftmost derivation}{no derivation}{a rightmost derivation in reverse}
\end{question}

\begin{question}{M}
In the table-driven predictive (LL(1)) parser, if the top of the stack is a non-terminal A and the current input symbol is a, the parser
\optslong{reduces by A}{reports success}{shifts a}{replaces A by the right side of the production in M[A, a], pushed in reverse}
\end{question}

\begin{question}{M}
In LL(1), the first L, the second L and the 1 stand for
\optslong{left-to-right scan, LR items, one table}{left-to-right scan, leftmost derivation, one symbol of lookahead}{leftmost scan, last derivation, one state}{left-recursive, left-factored, one production}
\end{question}

\begin{question}{M}
For the grammar S $\rightarrow$ ABC, A $\rightarrow$ a | $\varepsilon$, B $\rightarrow$ b | $\varepsilon$, C $\rightarrow$ c, the set FIRST(S) is
\optsshort{\{a, b, c\}}{\{a\}}{\{a, b\}}{\{a, b, c, $\varepsilon$\}}
\end{question}

\begin{question}{M}
For the grammar S $\rightarrow$ aSb | $\varepsilon$, the set FOLLOW(S) is
\optsshort{\{b\}}{\{a, b, \$\}}{\{\$\}}{\{b, \$\}}
\end{question}

\begin{question}{M}
For the grammar S $\rightarrow$ AB, A $\rightarrow$ a | $\varepsilon$, B $\rightarrow$ b | $\varepsilon$, the set FOLLOW(A) is
\optsshort{\{b, \$\}}{\{\$\}}{\{a, b, \$\}}{\{b\}}
\end{question}

\begin{question}{M}
Given below are two statements.\\
Statement I: Every unambiguous grammar is LL(1).\\
Statement II: Every LL(1) grammar is unambiguous.
\optslong{Both Statement I and Statement II are true}{Both Statement I and Statement II are false}{Statement I is true but Statement II is false}{Statement I is false but Statement II is true}
\end{question}

\begin{question}{M}
Panic-mode error recovery in an LL(1) parser uses
\optslong{backtracking}{the synchronizing set of tokens, such as FOLLOW(A), to skip input until a token that can follow A is found}{the LR automaton}{reduction by the handle}
\end{question}

\begin{question}{M}
A grammar G has productions S $\rightarrow$ aSb | ab. Which transformation makes it LL(1)?
\optslong{Removing left recursion}{Left factoring: S $\rightarrow$ aS', S' $\rightarrow$ Sb | b}{Converting to CNF}{Adding $\varepsilon$-productions}
\end{question}

\begin{question}{M}
Given below are two statements: one is labelled as Assertion A and the other is labelled as Reason R.\\
Assertion A: LL(1) parsers are more powerful than LR(1) parsers.\\
Reason R: Every LL(1) grammar is also LR(1).
\optslong{Both A and R are true and R is the correct explanation of A}{Both A and R are true but R is NOT the correct explanation of A}{A is true but R is false}{A is false but R is true}
\end{question}

\begin{question}{M}
The number of columns in an LL(1) parsing table for a grammar with 5 terminals is
\optsshort{7}{10}{6}{5}
\end{question}

\begin{question}{L}
Top-down parsing begins with
\optslong{a reduction}{the start symbol of the grammar}{the last token}{the input string}
\end{question}

\begin{question}{M}
Which of the following can cause backtracking in a top-down parser that is not predictive?
\optslong{Unique FIRST sets}{Productions of a non-terminal whose alternatives begin with the same terminal}{An LL(1) grammar}{An unambiguous grammar}
\end{question}

\begin{question}{M}
After eliminating left recursion from E $\rightarrow$ E + T | T, the result is
\optsshort{E $\rightarrow$ T + E | T}{E $\rightarrow$ E'T, E' $\rightarrow$ E + | $\varepsilon$}{E $\rightarrow$ +TE | T}{E $\rightarrow$ TE', E' $\rightarrow$ +TE' | $\varepsilon$}
\end{question}

\begin{question}{M}
Given below are two statements.\\
Statement I: Indirect left recursion, such as S $\rightarrow$ Aa, A $\rightarrow$ Sb | c, cannot be removed.\\
Statement II: FIRST($\alpha$) contains $\varepsilon$ exactly when $\alpha$ can derive the empty string.
\optslong{Both Statement I and Statement II are true}{Both Statement I and Statement II are false}{Statement I is true but Statement II is false}{Statement I is false but Statement II is true}
\end{question}

\begin{answers}
\ans{1}{2}{E' can start with + or derive $\varepsilon$.}
\ans{2}{1}{FOLLOW(T) contains FIRST(E') $-$ \{$\varepsilon$\} = \{+\} and, because E' can vanish, FOLLOW(E') = \{), \$\}.}
\ans{3}{3}{FOLLOW(F) includes FIRST(T') $-$ \{$\varepsilon$\} = \{*\} and FOLLOW(T) = \{+, ), \$\}.}
\ans{4}{4}{\$ because E is the start symbol, and ) from F $\rightarrow$ (E).}
\ans{5}{3}{An $\varepsilon$-production is placed under every terminal in FOLLOW(E') = \{), \$\}.}
\ans{6}{1}{Then each table entry holds at most one production.}
\ans{7}{1}{Left recursion makes FIRST sets of the alternatives overlap, and ambiguity gives two entries in some cell.}
\ans{8}{3}{The grammar is ambiguous; in practice the conflict is resolved by choosing S' $\rightarrow$ eS.}
\ans{9}{2}{FIRST(AaAb) = \{a\} and FIRST(BbBa) = \{b\} are disjoint, but in SLR the reductions A $\rightarrow$ $\varepsilon$ and B $\rightarrow$ $\varepsilon$ both have lookahead set \{a, b\}.}
\ans{10}{2}{The other grammars are left-recursive, need left factoring, or are ambiguous.}
\ans{11}{4}{Both alternatives begin with a.}
\ans{12}{4}{Without backtracking it is a predictive parser for LL(1) grammars.}
\ans{13}{1}{Left recursion must be eliminated before top-down parsing.}
\ans{14}{2}{A bottom-up (LR) parser produces a rightmost derivation in reverse.}
\ans{15}{4}{If M[A, a] is empty, it reports a syntax error.}
\ans{16}{2}{LR(k) means left-to-right scan, rightmost derivation in reverse, k symbols of lookahead.}
\ans{17}{1}{A and B can both vanish, so c can also begin S; C cannot vanish, so $\varepsilon$ is not included.}
\ans{18}{4}{\$ because S is the start symbol, and b from S $\rightarrow$ aSb.}
\ans{19}{1}{FIRST(B) = \{b, $\varepsilon$\}; since B can vanish, FOLLOW(S) = \{\$\} is added.}
\ans{20}{4}{For example, unambiguous left-recursive grammars are not LL(1).}
\ans{21}{2}{The non-terminal is popped and parsing resumes.}
\ans{22}{2}{Both alternatives start with a.}
\ans{23}{4}{R is true, and it shows that LR(1) is the more powerful class, so A is false.}
\ans{24}{3}{One column per terminal, plus one for the end marker \$.}
\ans{25}{2}{It expands non-terminals until the input string is produced.}
\ans{26}{2}{The parser may have to try the alternatives one by one.}
\ans{27}{4}{A $\rightarrow$ A$\alpha$ | $\beta$ becomes A $\rightarrow$ $\beta$A', A' $\rightarrow$ $\alpha$A' | $\varepsilon$.}
\ans{28}{4}{Indirect left recursion is removed by substituting productions in a fixed order of non-terminals.}
\end{answers}

\section{Syntax Analysis: Bottom-Up Parsing (shift-reduce, operator precedence, LR(0), SLR, CLR, LALR)}

\begin{question}{H}
The number of states in the canonical collection of LR(0) items for the augmented expression grammar E' $\rightarrow$ E, E $\rightarrow$ E + T | T, T $\rightarrow$ T * F | F, F $\rightarrow$ (E) | id is
\optsshort{9}{10}{12}{14}
\end{question}

\begin{question}{H}
For the grammar S $\rightarrow$ CC, C $\rightarrow$ cC | d, the numbers of states in the LR(0) (SLR), canonical LR(1) and LALR(1) automata are respectively
\optsshort{10, 10 and 7}{7, 7 and 10}{7, 10 and 7}{7, 10 and 10}
\end{question}

\begin{question}{H}
Match List I with List II.\\
List I (Parser): A. LR(0)  B. SLR(1)  C. LALR(1)  D. Canonical LR(1)\\
List II (Lookahead used for reductions): I. LR(1) items, with states having identical cores merged  II. No lookahead: reduce whenever a complete item is present  III. FOLLOW set of the left-hand non-terminal  IV. Exact lookahead symbols carried in LR(1) items
\optsshort{A-II, B-III, C-IV, D-I}{A-II, B-III, C-I, D-IV}{A-II, B-I, C-III, D-IV}{A-III, B-II, C-I, D-IV}
\end{question}

\begin{question}{H}
Arrange the following parsers in increasing order of the class of grammars they can handle.\\
A. LALR(1)  B. LR(0)  C. Canonical LR(1)  D. SLR(1)
\optsshort{B, D, C, A}{B, D, A, C}{B, A, D, C}{D, B, A, C}
\end{question}

\begin{question}{H}
Given below are two statements.\\
Statement I: The grammar S $\rightarrow$ L = R | R, L $\rightarrow$ *R | id, R $\rightarrow$ L is not SLR(1).\\
Statement II: The grammar S $\rightarrow$ L = R | R, L $\rightarrow$ *R | id, R $\rightarrow$ L is LALR(1).
\optslong{Both Statement I and Statement II are true}{Both Statement I and Statement II are false}{Statement I is true but Statement II is false}{Statement I is false but Statement II is true}
\end{question}

\begin{question}{H}
Given below are two statements: one is labelled as Assertion A and the other is labelled as Reason R.\\
Assertion A: Merging LR(1) states with the same core to form LALR(1) states can introduce reduce-reduce conflicts but never shift-reduce conflicts.\\
Reason R: The shift actions of a state depend only on its core, so a shift-reduce conflict after merging would already exist in one of the original LR(1) states.
\optslong{Both A and R are true and R is the correct explanation of A}{Both A and R are true but R is NOT the correct explanation of A}{A is true but R is false}{A is false but R is true}
\end{question}

\begin{question}{H}
The grammar S $\rightarrow$ Aa | bAc | dc | bda, A $\rightarrow$ d is
\optsshort{neither SLR(1) nor LALR(1)}{SLR(1) and LALR(1)}{LALR(1) but not SLR(1)}{SLR(1) but not LALR(1)}
\end{question}

\begin{question}{H}
The grammar S $\rightarrow$ aAd | bBd | aBe | bAe, A $\rightarrow$ c, B $\rightarrow$ c is
\optsshort{LR(1) but not LALR(1)}{SLR(1)}{not LR(1)}{LALR(1)}
\end{question}

\begin{question}{H}
In shift-reduce parsing, a handle is
\optslong{the first symbol on the stack}{a substring that matches the right side of a production and whose reduction is one step of a rightmost derivation in reverse}{any substring of the right sentential form}{the lookahead symbol}
\end{question}

\begin{question}{H}
Match List I with List II.\\
List I (Conflict or term): A. Shift-reduce conflict  B. Reduce-reduce conflict  C. Viable prefix  D. Handle pruning\\
List II (Meaning): I. A prefix of a right sentential form that can appear on the stack of a shift-reduce parser  II. The parser cannot decide whether to shift or to reduce  III. The parser cannot decide which of two productions to reduce by  IV. Repeatedly reducing handles to obtain the rightmost derivation in reverse
\optsshort{A-II, B-I, C-III, D-IV}{A-II, B-III, C-IV, D-I}{A-III, B-II, C-I, D-IV}{A-II, B-III, C-I, D-IV}
\end{question}

\begin{question}{H}
The four actions of an LR parser are
\optslong{shift, goto, follow and first}{scan, parse, check and emit}{push, pop, read and write}{shift, reduce, accept and error}
\end{question}

\begin{question}{H}
Given below are two statements.\\
Statement I: An ambiguous grammar can never be LR(k) for any k.\\
Statement II: Yacc can use an ambiguous grammar if precedence and associativity declarations resolve the conflicts.
\optslong{Both Statement I and Statement II are true}{Both Statement I and Statement II are false}{Statement I is true but Statement II is false}{Statement I is false but Statement II is true}
\end{question}

\begin{question}{M}
In an LR(0) automaton, a state containing both A $\rightarrow$ $\alpha\cdot$ and B $\rightarrow$ $\beta\cdot$a$\gamma$ has a
\optsshort{no conflict}{goto conflict}{shift-reduce conflict}{reduce-reduce conflict}
\end{question}

\begin{question}{M}
For the expression grammar E $\rightarrow$ E + T | T, T $\rightarrow$ T * F | F, F $\rightarrow$ (E) | id, which statement is correct?
\optsshort{It is LR(0)}{It is LL(1)}{It is SLR(1) but not LR(0)}{It is ambiguous}
\end{question}

\begin{question}{M}
An operator grammar is one in which
\optslong{the grammar is in CNF}{all productions are left-recursive}{no production has $\varepsilon$ on its right side and no right side has two adjacent non-terminals}{every production has an operator}
\end{question}

\begin{question}{M}
In operator-precedence parsing, the relation a <$\cdot$ b between two terminals means
\optslong{a and b have equal precedence}{a has higher precedence than b}{a yields precedence to b (b's handle starts after a)}{a and b cannot be adjacent}
\end{question}

\begin{question}{M}
Given below are two statements: one is labelled as Assertion A and the other is labelled as Reason R.\\
Assertion A: Canonical LR(1) parsers have many more states than LALR(1) parsers for typical programming languages.\\
Reason R: LALR(1) parsers have the same number of states as SLR(1) parsers.
\optslong{Both A and R are true and R is the correct explanation of A}{Both A and R are true but R is NOT the correct explanation of A}{A is true but R is false}{A is false but R is true}
\end{question}

\begin{question}{M}
Which parser generator tool produces LALR(1) parsers?
\optsshort{Lex}{Yacc (and Bison)}{ANTLR (LL(*))}{Flex}
\end{question}

\begin{question}{M}
An LR(0) item is
\optslong{a production with a dot at some position on its right side}{a terminal}{a parse tree}{a token}
\end{question}

\begin{question}{M}
An LR(1) item has the form
\optslong{[A $\rightarrow$ $\alpha\cdot\beta$]}{[$\alpha$, $\beta$, a, b]}{[A $\rightarrow$ $\alpha\cdot\beta$, a], where a is a lookahead terminal}{[A, a]}
\end{question}

\begin{question}{M}
Given below are two statements.\\
Statement I: Every SLR(1) grammar is LALR(1).\\
Statement II: Every LALR(1) grammar is SLR(1).
\optslong{Both Statement I and Statement II are true}{Both Statement I and Statement II are false}{Statement I is true but Statement II is false}{Statement I is false but Statement II is true}
\end{question}

\begin{question}{M}
A bottom-up parser builds the parse tree
\optsshort{it does not build a tree}{in level order}{from the leaves to the root}{from the root to the leaves}
\end{question}

\begin{question}{M}
For the grammar S $\rightarrow$ aSb | ab and input aabb, the sequence of actions of a shift-reduce parser is
\optslong{shift, reduce, shift, reduce, accept}{reduce, shift, shift, shift, accept}{shift, shift, shift, reduce S $\rightarrow$ ab, shift, reduce S $\rightarrow$ aSb, accept}{shift, shift, reduce, reduce, accept}
\end{question}

\begin{question}{M}
The LR parsing method is preferred in practice because
\optslong{it needs no tables}{it is easy to write by hand}{it can handle virtually all programming-language constructs that have CFGs, detects errors as early as possible, and is more general than LL}{it works only for regular languages}
\end{question}

\begin{question}{M}
Given below are two statements: one is labelled as Assertion A and the other is labelled as Reason R.\\
Assertion A: The LALR(1) parser for a grammar has the same number of states as its SLR(1) parser.\\
Reason R: Yacc generates LALR(1) parsers.
\optslong{Both A and R are true and R is the correct explanation of A}{Both A and R are true but R is NOT the correct explanation of A}{A is true but R is false}{A is false but R is true}
\end{question}

\begin{question}{M}
For the grammar S $\rightarrow$ CC, C $\rightarrow$ cC | d, the LALR(1) parser is
\optslong{not constructible}{larger than the canonical LR(1) parser}{identical in size to the canonical LR(1) parser}{smaller than the canonical LR(1) parser and has no conflicts}
\end{question}

\begin{question}{M}
Given below are two statements.\\
Statement I: An LR(0) parser can parse any unambiguous grammar.\\
Statement II: An LR(0) grammar cannot have an $\varepsilon$-production whose reduction conflicts with a shift.
\optslong{Both Statement I and Statement II are true}{Both Statement I and Statement II are false}{Statement I is true but Statement II is false}{Statement I is false but Statement II is true}
\end{question}

\begin{question}{L}
In LR(k), the letter R indicates that the parser constructs
\optslong{a reduced grammar}{a regular expression}{a recursive derivation}{a rightmost derivation in reverse}
\end{question}

\begin{question}{M}
In an SLR parser, the action reduce by A $\rightarrow$ $\alpha$ is placed in state I for terminal a when A $\rightarrow$ $\alpha\cdot$ is in I and
\optslong{a is in FOLLOW(A)}{a is the next token in the input always}{a is in FIRST(A)}{a is in FIRST($\alpha$)}
\end{question}

\begin{question}{M}
The dangling-else conflict in an LR parser is a
\optslong{lexical conflict}{shift-reduce conflict, usually resolved by shifting}{reduce-reduce conflict}{goto conflict}
\end{question}

\begin{answers}
\ans{1}{3}{States I0 to I11, as in the Dragon Book.}
\ans{2}{3}{LALR merges LR(1) states with the same core: 10 states become 7.}
\ans{3}{2}{Power increases from LR(0) to SLR, LALR and canonical LR.}
\ans{4}{2}{LR(0) $\subset$ SLR(1) $\subset$ LALR(1) $\subset$ LR(1).}
\ans{5}{1}{In SLR, = is in FOLLOW(R), causing a shift-reduce conflict between shifting = and reducing R $\rightarrow$ L; LALR lookaheads avoid it.}
\ans{6}{1}{The grammar S $\rightarrow$ aAd | bBd | aBe | bAe, A $\rightarrow$ c, B $\rightarrow$ c is LR(1) but not LALR(1) because of reduce-reduce conflicts.}
\ans{7}{3}{FOLLOW(A) = \{a, c\}, so SLR has a conflict on a after d, while LALR lookaheads distinguish the cases.}
\ans{8}{1}{The LR(1) states for A $\rightarrow$ c$\cdot$ and B $\rightarrow$ c$\cdot$ have the same core and merge into a state with reduce-reduce conflicts.}
\ans{9}{2}{Reducing a handle moves the parse one step back towards the start symbol.}
\ans{10}{4}{Basic shift-reduce terminology.}
\ans{11}{4}{The GOTO table is used after a reduction to find the next state.}
\ans{12}{1}{Yacc's default is to prefer shift in a shift-reduce conflict and the earlier production in a reduce-reduce conflict.}
\ans{13}{3}{The parser could reduce by A $\rightarrow$ $\alpha$ or shift a.}
\ans{14}{3}{Some LR(0) states contain both a complete item and a shift item (e.g. E $\rightarrow$ T$\cdot$ and T $\rightarrow$ T$\cdot$*F), but FOLLOW sets resolve them.}
\ans{15}{3}{Operator-precedence parsing works on such grammars.}
\ans{16}{3}{a $\cdot$> b means a takes precedence; a $\doteq$ b means they are in the same handle.}
\ans{17}{2}{Both are true, but R does not explain A; the LR(1) blow-up comes from splitting states by lookahead.}
\ans{18}{2}{Yacc stands for "Yet Another Compiler-Compiler".}
\ans{19}{1}{For example A $\rightarrow$ X$\cdot$YZ means X has been seen and YZ is expected.}
\ans{20}{3}{The lookahead is used only when the dot is at the right end.}
\ans{21}{3}{S $\rightarrow$ L = R | R, L $\rightarrow$ *R | id, R $\rightarrow$ L is LALR(1) but not SLR(1).}
\ans{22}{3}{It reduces the input to the start symbol.}
\ans{23}{3}{After shifting a, a, b the handle ab is on top; reduce it, shift b and reduce aSb.}
\ans{24}{3}{LR parsers are usually generated by tools.}
\ans{25}{2}{Both are true, but A holds because both parsers are built from the LR(0) automaton.}
\ans{26}{4}{7 states instead of 10, with no conflicts.}
\ans{27}{4}{Many unambiguous grammars, such as the standard expression grammar, are not LR(0).}
\ans{28}{4}{L stands for left-to-right scanning of the input.}
\ans{29}{1}{This use of FOLLOW sets is what can cause spurious conflicts.}
\ans{30}{2}{Shifting associates else with the nearest unmatched if.}
\end{answers}

\section{Semantic Analysis (attribute grammars, syntax-directed definitions and translation, type checking)}

\begin{question}{H}
Consider the translation scheme E $\rightarrow$ E1 \# T \{E.val = E1.val $\times$ T.val\} | T \{E.val = T.val\}, T $\rightarrow$ T1 \& F \{T.val = T1.val + F.val\} | F \{T.val = F.val\}, F $\rightarrow$ num \{F.val = num.val\}. The value computed for 2 \# 3 \& 5 \# 6 \& 4 is
\optsshort{200}{160}{40}{72}
\end{question}

\begin{question}{H}
The translation scheme E $\rightarrow$ E + T \{print("+")\} | T, T $\rightarrow$ num \{print(num.val)\} prints what for the input 2 + 3 + 4?
\optsshort{23+4+}{2+3+4}{+2+34}{234++}
\end{question}

\begin{question}{H}
Match List I with List II.\\
List I (Concept): A. Synthesized attribute  B. Inherited attribute  C. S-attributed definition  D. L-attributed definition\\
List II (Meaning): I. Uses only synthesized attributes  II. Computed from the attributes of the node's children  III. Inherited attributes depend only on the parent and on siblings to the left  IV. Computed from the attributes of the parent and/or siblings
\optsshort{A-IV, B-II, C-I, D-III}{A-II, B-IV, C-III, D-I}{A-II, B-IV, C-I, D-III}{A-II, B-I, C-IV, D-III}
\end{question}

\begin{question}{H}
Given below are two statements.\\
Statement I: Every S-attributed definition is L-attributed.\\
Statement II: Every L-attributed definition is S-attributed.
\optslong{Both Statement I and Statement II are true}{Both Statement I and Statement II are false}{Statement I is true but Statement II is false}{Statement I is false but Statement II is true}
\end{question}

\begin{question}{H}
S-attributed definitions are naturally evaluated during
\optslong{bottom-up (LR) parsing, by computing attributes when reducing}{code optimization}{lexical analysis}{top-down parsing using inherited attributes}
\end{question}

\begin{question}{H}
Given below are two statements: one is labelled as Assertion A and the other is labelled as Reason R.\\
Assertion A: L-attributed definitions can be evaluated in a single depth-first, left-to-right traversal of the parse tree.\\
Reason R: Every inherited attribute depends only on information from the parent and from siblings to its left, which is already available.
\optslong{Both A and R are true and R is the correct explanation of A}{Both A and R are true but R is NOT the correct explanation of A}{A is true but R is false}{A is false but R is true}
\end{question}

\begin{question}{H}
An attribute grammar has a dependency graph with a cycle. This means that
\optslong{evaluation is trivial}{there is no order in which all attributes can be evaluated}{the grammar is S-attributed}{the grammar is ambiguous}
\end{question}

\begin{question}{M}
In the declaration grammar D $\rightarrow$ T L, T $\rightarrow$ int | float, L $\rightarrow$ L1 , id | id, the type is passed from T to the identifiers in L using
\optslong{a global variable}{a synthesized attribute only}{the symbol table only}{an inherited attribute L.inh = T.type}
\end{question}

\begin{question}{M}
A parse tree showing the values of attributes at each node is called
\optslong{a syntax tree}{an activation tree}{a DAG}{an annotated (decorated) parse tree}
\end{question}

\begin{question}{M}
Given below are two statements.\\
Statement I: A syntax-directed definition specifies semantic rules without specifying when they are evaluated.\\
Statement II: A syntax-directed translation scheme embeds program fragments (semantic actions) within production bodies, fixing the evaluation order.
\optslong{Both Statement I and Statement II are true}{Both Statement I and Statement II are false}{Statement I is true but Statement II is false}{Statement I is false but Statement II is true}
\end{question}

\begin{question}{H}
Match List I with List II.\\
List I (Type concept): A. Coercion  B. Cast  C. Overloading  D. Parametric polymorphism\\
List II (Description): I. One operator or function name with several meanings depending on the argument types  II. Explicit type conversion written by the programmer  III. Implicit type conversion inserted by the compiler  IV. A single definition that works uniformly for many types, such as a generic list
\optsshort{A-III, B-II, C-IV, D-I}{A-II, B-III, C-I, D-IV}{A-III, B-I, C-II, D-IV}{A-III, B-II, C-I, D-IV}
\end{question}

\begin{question}{H}
Given below are two statements.\\
Statement I: Static type checking is done at compile time, while dynamic type checking is done at run time.\\
Statement II: A strongly typed language guarantees that a program accepted by its compiler will run without type errors.
\optslong{Both Statement I and Statement II are true}{Both Statement I and Statement II are false}{Statement I is true but Statement II is false}{Statement I is false but Statement II is true}
\end{question}

\begin{question}{M}
Under structural equivalence of types, two type expressions are equivalent if
\optslong{they have the same name}{they are declared in the same scope}{they are built from the same basic types by the same type constructors}{they have the same size}
\end{question}

\begin{question}{M}
In C, the declarations "typedef int A; typedef int B;" are compatible with each other because C uses
\optslong{no type checking}{dynamic typing}{name equivalence for typedefs}{structural (or rather typedef-transparent) equivalence, since typedef only creates an alias}
\end{question}

\begin{question}{M}
Type inference in languages such as ML and Haskell is performed by
\optslong{lexical analysis}{guessing types at run time}{ignoring types}{unification of type expressions (Hindley-Milner algorithm)}
\end{question}

\begin{question}{M}
Given below are two statements: one is labelled as Assertion A and the other is labelled as Reason R.\\
Assertion A: Checking that a variable is declared before use cannot be done by a context-free grammar alone.\\
Reason R: The language \{wcw : w is a string\} models declaration before use, and it is not context-free.
\optslong{Both A and R are true and R is the correct explanation of A}{Both A and R are true but R is NOT the correct explanation of A}{A is true but R is false}{A is false but R is true}
\end{question}

\begin{question}{M}
The type expression for an array of 10 pointers to integers is
\optslong{function(integer, pointer)}{pointer(array(0..9, integer))}{integer $\times$ 10}{array(0..9, pointer(integer))}
\end{question}

\begin{question}{M}
For the expression x + y, where x is an int and y is a float, a C compiler
\optslong{reports a type error}{converts x to float (widening) and performs floating-point addition}{uses integer addition}{converts y to int}
\end{question}

\begin{question}{M}
In an SDT, a semantic action placed at the right end of a production is executed
\optslong{never}{before the production is applied}{when the body has been completely recognised (on reduction in a bottom-up parser)}{only during lexical analysis}
\end{question}

\begin{question}{M}
Given below are two statements.\\
Statement I: Inherited attributes are useful for passing context information, such as the expected type, down the parse tree.\\
Statement II: Synthesized attributes of terminals are supplied by the lexical analyzer.
\optslong{Both Statement I and Statement II are true}{Both Statement I and Statement II are false}{Statement I is true but Statement II is false}{Statement I is false but Statement II is true}
\end{question}

\begin{question}{M}
For the SDD with productions S $\rightarrow$ AB \{A.i = 1; B.i = A.s; S.s = B.s\}, A $\rightarrow$ a \{A.s = A.i + 1\}, B $\rightarrow$ b \{B.s = B.i $\times$ 2\}, the value of S.s for the input ab is
\optsshort{3}{6}{2}{4}
\end{question}

\begin{question}{M}
Which kind of SDD can always be evaluated in any bottom-up parser during parsing?
\optslong{L-attributed definitions}{Definitions with cyclic dependencies}{S-attributed definitions}{Any SDD}
\end{question}

\begin{question}{L}
The semantic analyzer uses which data structure to check that identifiers are declared and used with the correct types?
\optsshort{Flow graph}{Activation tree}{Symbol table}{Parse table}
\end{question}

\begin{question}{M}
Given below are two statements: one is labelled as Assertion A and the other is labelled as Reason R.\\
Assertion A: Every L-attributed definition can be evaluated during LR parsing without any modification.\\
Reason R: LR parsers know the production being used before they reach its right end.
\optslong{Both A and R are true and R is the correct explanation of A}{Both A and R are true but R is NOT the correct explanation of A}{A is true but R is false}{A is false but R is true}
\end{question}

\begin{question}{M}
A dependency graph for attribute evaluation has an edge from attribute b to attribute c when
\optslong{c is computed using b}{c is inherited}{b and c belong to the same node}{b and c are both synthesized}
\end{question}

\begin{answers}
\ans{1}{2}{\& binds tighter than \# (it is lower in the grammar), so the value is 2 $\times$ (3 + 5) $\times$ (6 + 4) = 160.}
\ans{2}{1}{The actions print the postfix form; + is left-associative.}
\ans{3}{3}{Standard attribute grammar terminology.}
\ans{4}{3}{L-attributed definitions may also have inherited attributes.}
\ans{5}{1}{Attribute values can be kept on the parser stack alongside grammar symbols.}
\ans{6}{1}{This is why L-attributed definitions suit top-down (LL) translation.}
\ans{7}{2}{Attribute evaluation follows a topological sort of an acyclic dependency graph.}
\ans{8}{4}{L1.inh = L.inh then passes the type down the list.}
\ans{9}{4}{Attribute values are written next to the nodes.}
\ans{10}{1}{This is the distinction between SDDs and SDTs.}
\ans{11}{4}{Type concepts handled during semantic analysis.}
\ans{12}{1}{Both are standard definitions.}
\ans{13}{3}{Under name equivalence, types are equivalent only if they have the same name.}
\ans{14}{4}{Structure types in C, by contrast, are compared by name (tag).}
\ans{15}{4}{The most general type of each expression is computed without declarations.}
\ans{16}{1}{Such checks are done during semantic analysis, using the symbol table.}
\ans{17}{4}{Type constructors are applied from the outside in.}
\ans{18}{2}{Implicit widening coercion follows the usual arithmetic conversions.}
\ans{19}{3}{Such postfix SDTs can be implemented directly during LR parsing.}
\ans{20}{1}{For example num.val comes from the token's attribute value.}
\ans{21}{4}{A.s = 1 + 1 = 2; B.i = 2; B.s = 4; S.s = 4.}
\ans{22}{3}{They need only synthesized attributes, computed at reduction time.}
\ans{23}{3}{Each identifier's entry records its type, scope and other attributes.}
\ans{24}{4}{An LR parser does not know the production until it reduces; marker non-terminals are needed to handle inherited attributes, so A is false.}
\ans{25}{1}{A topological sort of the graph gives a valid evaluation order.}
\end{answers}

\section{Run-Time Environments (storage organization, activation records, parameter passing, scoping, symbol tables)}

\begin{question}{H}
Consider: procedure p(x, y) \{ x = x + 1; y = y + 1; \} and the main program a = 1; p(a, a); print(a). The values printed under call by value, call by reference and call by value-result are respectively
\optsshort{3, 3 and 2}{1, 2 and 3}{1, 3 and 3}{1, 3 and 2}
\end{question}

\begin{question}{H}
Consider swap(x, y) \{ t = x; x = y; y = t; \} called as swap(a, b) with a = 1 and b = 2. The values of a and b afterwards under call by value and under call by reference are respectively
\optsshort{2, 1 and 2, 1}{1, 2 and 1, 2}{2, 1 and 1, 2}{1, 2 and 2, 1}
\end{question}

\begin{question}{H}
Consider: int x = 1; void f() \{ print(x); \} void g() \{ int x = 2; f(); \} and main calls g(). The values printed under static and dynamic scoping are respectively
\optsshort{2 and 1}{1 and 1}{1 and 2}{2 and 2}
\end{question}

\begin{question}{H}
Match List I with List II.\\
List I (Parameter passing method): A. Call by value  B. Call by reference  C. Call by value-result  D. Call by name\\
List II (Description): I. The address of the actual parameter is passed  II. The actual parameter is textually substituted and re-evaluated on each use  III. The value is copied in, and the final value is copied back on return  IV. The value of the actual parameter is copied into the formal
\optsshort{A-IV, B-III, C-I, D-II}{A-IV, B-I, C-III, D-II}{A-I, B-IV, C-III, D-II}{A-IV, B-I, C-II, D-III}
\end{question}

\begin{question}{H}
Match List I with List II.\\
List I (Activation record field): A. Control link  B. Access link  C. Saved machine status  D. Return value\\
List II (Purpose): I. Points to the activation record of the lexically enclosing procedure, for non-local data  II. Points to the activation record of the caller  III. Holds the return address and register contents  IV. Space for the value returned to the caller
\optsshort{A-I, B-II, C-III, D-IV}{A-II, B-I, C-IV, D-III}{A-II, B-III, C-I, D-IV}{A-II, B-I, C-III, D-IV}
\end{question}

\begin{question}{H}
Given below are two statements.\\
Statement I: Static allocation of activation records cannot support recursion.\\
Statement II: Stack allocation of activation records supports recursion.
\optslong{Both Statement I and Statement II are true}{Both Statement I and Statement II are false}{Statement I is true but Statement II is false}{Statement I is false but Statement II is true}
\end{question}

\begin{question}{H}
main calls fact(5), where fact(n) returns 1 if n $\leq$ 1 and n $\times$ fact(n $-$ 1) otherwise. The maximum number of activation records on the control stack, including main's, is
\optsshort{7}{6}{5}{10}
\end{question}

\begin{question}{H}
In the typical run-time memory layout of a C program,
\optslong{there is no heap}{the code and static data are at fixed locations, and the heap and stack grow towards each other}{the heap is inside the stack}{the stack and the heap both grow towards higher addresses}
\end{question}

\begin{question}{M}
An activation tree represents
\optslong{the flow graph}{the parse tree}{the sequence of procedure activations during one run, with each child a call made by its parent}{the call graph of a program at compile time}
\end{question}

\begin{question}{H}
Given below are two statements: one is labelled as Assertion A and the other is labelled as Reason R.\\
Assertion A: Reference counting reclaims cyclic garbage as soon as the cycle becomes unreachable.\\
Reason R: Objects in an unreachable cycle still have non-zero reference counts.
\optslong{Both A and R are true and R is the correct explanation of A}{Both A and R are true but R is NOT the correct explanation of A}{A is true but R is false}{A is false but R is true}
\end{question}

\begin{question}{H}
Match List I with List II.\\
List I (Garbage collection technique): A. Reference counting  B. Mark-and-sweep  C. Copying collection  D. Generational collection\\
List II (Characteristic): I. Moves live objects to a new space, compacting memory  II. Collects young objects more often, since most objects die young  III. Reclaims an object as soon as its count drops to zero  IV. Marks reachable objects from the roots and then frees the unmarked ones
\optsshort{A-III, B-IV, C-I, D-II}{A-III, B-IV, C-II, D-I}{A-III, B-I, C-IV, D-II}{A-IV, B-III, C-I, D-II}
\end{question}

\begin{question}{M}
A dangling reference occurs when
\optslong{a pointer refers to storage that has already been deallocated}{memory is never freed}{the stack overflows}{a variable is not initialised}
\end{question}

\begin{question}{M}
A display is used to
\optslong{allocate heap memory}{access non-local variables in nested procedures quickly through an array of pointers to activation records at each nesting depth}{show the output of a program}{store the symbol table}
\end{question}

\begin{question}{M}
In a language with nested procedures, accessing a variable declared at nesting depth 1 from a procedure at nesting depth 3 using access links requires following
\optsshort{2 access links}{3 access links}{1 access link}{no links}
\end{question}

\begin{question}{M}
Given below are two statements.\\
Statement I: C passes an array parameter by copying all of its elements.\\
Statement II: In Java, assigning a new object to a parameter inside a method changes the caller's variable.
\optslong{Both Statement I and Statement II are true}{Both Statement I and Statement II are false}{Statement I is true but Statement II is false}{Statement I is false but Statement II is true}
\end{question}

\begin{question}{M}
Which data structure is most commonly used to implement a symbol table in a compiler for fast look-up?
\optsshort{Stack}{Hash table}{Queue}{Linear list}
\end{question}

\begin{question}{M}
To handle nested block scopes, a compiler can
\optslong{use only static allocation}{use a single global table with no scope information}{maintain a separate symbol table for each scope, linked to the table of the enclosing scope}{ignore inner declarations}
\end{question}

\begin{question}{M}
Given below are two statements: one is labelled as Assertion A and the other is labelled as Reason R.\\
Assertion A: Local variables of a recursive procedure must be allocated on the stack (or heap), not statically.\\
Reason R: Several activations of the procedure may be alive at the same time, each needing its own copy of the locals.
\optslong{Both A and R are true and R is the correct explanation of A}{Both A and R are true but R is NOT the correct explanation of A}{A is true but R is false}{A is false but R is true}
\end{question}

\begin{question}{M}
The calling sequence of a procedure call is code that
\optslong{optimizes the procedure body}{links object files}{allocates the activation record and fills in its fields on the caller and callee sides}{parses the procedure}
\end{question}

\begin{question}{M}
Under call by name, a procedure with formal parameter x evaluates the actual parameter expression
\optslong{only if x is assigned}{never}{each time x is used in the body}{once, before the call}
\end{question}

\begin{question}{M}
Heap allocation is needed for
\optslong{local variables of non-recursive procedures}{register variables}{global constants}{data whose lifetime is not tied to procedure activations, such as dynamically created objects}
\end{question}

\begin{question}{M}
Given below are two statements.\\
Statement I: The access link of a procedure's activation record points to its caller.\\
Statement II: In a language without nested procedures, such as C, access links are not needed.
\optslong{Both Statement I and Statement II are true}{Both Statement I and Statement II are false}{Statement I is true but Statement II is false}{Statement I is false but Statement II is true}
\end{question}

\begin{question}{L}
The area of memory used to hold the activation records of procedures during execution is the
\optsshort{code area}{heap}{control stack}{static area}
\end{question}

\begin{question}{M}
A function in C returns a pointer to one of its local (automatic) variables. This is an error because
\optslong{local variables are stored in the heap}{the local variable is freed when the function returns, leaving a dangling pointer}{C does not allow pointers}{the pointer is copied by value}
\end{question}

\begin{answers}
\ans{1}{4}{By value: a is unchanged. By reference: x and y are both a, incremented twice. By value-result: x and y become 2 and are copied back to a.}
\ans{2}{4}{Call by value swaps only the local copies.}
\ans{3}{3}{Static scoping uses the x in f's defining environment (global); dynamic scoping uses the most recent active binding (g's x).}
\ans{4}{2}{Call by name is implemented with thunks (as in ALGOL 60).}
\ans{5}{4}{The control link is also called the dynamic link, and the access link the static link.}
\ans{6}{1}{With static allocation (as in early FORTRAN) each procedure has only one activation record.}
\ans{7}{2}{main, fact(5), fact(4), fact(3), fact(2) and fact(1) are active together.}
\ans{8}{2}{The heap usually grows upward and the stack downward.}
\ans{9}{3}{The control stack at any time holds the activations on a path from the root.}
\ans{10}{4}{Because of R, reference counting cannot reclaim cycles; tracing collectors can.}
\ans{11}{1}{Common garbage-collection strategies.}
\ans{12}{1}{For example, returning the address of a local variable in C.}
\ans{13}{2}{It avoids following a chain of access links.}
\ans{14}{1}{The number of links equals the difference in nesting depths.}
\ans{15}{2}{In C the array name decays to a pointer; Java passes references by value, so reassigning the parameter has no effect on the caller.}
\ans{16}{2}{Scoping is handled by a stack of tables or by chaining entries by scope.}
\ans{17}{3}{Look-up searches the innermost table first and then outward.}
\ans{18}{1}{Static allocation gives one copy per procedure.}
\ans{19}{3}{The return sequence restores the machine state and returns.}
\ans{20}{3}{Jensen's device exploits this behaviour.}
\ans{21}{4}{It is managed with malloc/free, new/delete or garbage collection.}
\ans{22}{4}{The access link points to the lexically enclosing procedure; the control link points to the caller.}
\ans{23}{3}{Activation records are pushed on calls and popped on returns.}
\ans{24}{2}{A static local variable would not have this problem.}
\end{answers}

\section{Intermediate Code Generation (syntax trees, DAGs, three-address code, quadruples, triples, SSA, backpatching)}

\begin{question}{H}
The three-address code for a = b * $-$c + b * $-$c, generated without detecting common subexpressions, has how many instructions (including the final assignment)?
\optsshort{6}{7}{4}{5}
\end{question}

\begin{question}{H}
If common subexpressions are detected using a DAG, the three-address code for a = b * $-$c + b * $-$c can be reduced to how many instructions?
\optsshort{4}{6}{3}{5}
\end{question}

\begin{question}{H}
Match List I with List II.\\
List I (Representation): A. Quadruples  B. Triples  C. Indirect triples  D. Static single-assignment form\\
List II (Feature): I. Each variable is assigned exactly once; $\varphi$-functions merge values at join points  II. Fields op, arg1, arg2 and result; temporaries are named explicitly  III. A separate list of pointers to triples, so instructions can be reordered easily  IV. Results are referred to by the position of the instruction that computes them
\optsshort{A-IV, B-II, C-III, D-I}{A-II, B-IV, C-III, D-I}{A-II, B-IV, C-I, D-III}{A-II, B-III, C-IV, D-I}
\end{question}

\begin{question}{H}
Given below are two statements: one is labelled as Assertion A and the other is labelled as Reason R.\\
Assertion A: Moving an instruction is harder in a triple representation than in a quadruple representation.\\
Reason R: In triples, other instructions refer to a result by its position, so moving it requires changing all references.
\optslong{Both A and R are true and R is the correct explanation of A}{Both A and R are true but R is NOT the correct explanation of A}{A is true but R is false}{A is false but R is true}
\end{question}

\begin{question}{H}
The syntax tree for a + a * (b $-$ c) + (b $-$ c) * d has 13 nodes. The corresponding DAG has
\optsshort{11 nodes}{9 nodes}{10 nodes}{7 nodes}
\end{question}

\begin{question}{M}
The minimum number of temporary variables needed to compute (a + b) * (c + d) in three-address code, if temporaries may be reused, is
\optsshort{2}{3}{1}{4}
\end{question}

\begin{question}{H}
In static single-assignment form, the code x = a; x = x + 1; y = x * 2 becomes
\optslong{x = a; x = x + 1; y = x * 2}{x1 = a; x2 = x1 + 1; y1 = x2 * 2}{x1 = a; x1 = x1 + 1; y1 = x1 * 2}{x = a + 1; y = x * 2}
\end{question}

\begin{question}{M}
A $\varphi$-function in SSA form is placed
\optslong{at the start of every basic block}{only in loops without branches}{after every assignment}{at a join point where different definitions of a variable reach along different control-flow paths}
\end{question}

\begin{question}{H}
For an integer array A[10][20] (4 bytes per element, indices from 0) stored in row-major order with base address B, the address of A[i][j] is
\optsshort{B + (i + j) $\times$ 4}{B + (i $\times$ 20 + j) $\times$ 4}{B + (j $\times$ 10 + i) $\times$ 4}{B + (i $\times$ 10 + j) $\times$ 4}
\end{question}

\begin{question}{H}
Backpatching is used in intermediate code generation to
\optslong{allocate registers}{build the symbol table}{optimize loops}{fill in jump targets of conditional and unconditional jumps that are not yet known when the jumps are generated}
\end{question}

\begin{question}{M}
With short-circuit (jumping) code for the boolean expression a < b or c < d, if a < b is true,
\optslong{the expression evaluates to false}{c < d is still evaluated}{control jumps directly to the true exit without evaluating c < d}{an error occurs}
\end{question}

\begin{question}{M}
Given below are two statements.\\
Statement I: Three-address code has at most one operator on the right side of each instruction.\\
Statement II: Postfix notation is a form of intermediate representation.
\optslong{Both Statement I and Statement II are true}{Both Statement I and Statement II are false}{Statement I is true but Statement II is false}{Statement I is false but Statement II is true}
\end{question}

\begin{question}{M}
The postfix form of the expression a + b * c is
\optsshort{a+bc*}{ab+c*}{abc*+}{+a*bc}
\end{question}

\begin{question}{M}
A procedure call p(x1, x2, x3) is translated into three-address code as
\optslong{p = x1 + x2 + x3}{param x1; param x2; param x3; call p, 3}{goto p}{call p, 3}
\end{question}

\begin{question}{M}
Given below are two statements: one is labelled as Assertion A and the other is labelled as Reason R.\\
Assertion A: Generating intermediate code makes it easier to apply machine-independent optimizations.\\
Reason R: Intermediate code is independent of the target machine and simpler than the source program.
\optslong{Both A and R are true and R is the correct explanation of A}{Both A and R are true but R is NOT the correct explanation of A}{A is true but R is false}{A is false but R is true}
\end{question}

\begin{question}{M}
A switch statement with many dense case values is efficiently translated using
\optslong{recursion}{a jump table indexed by the case value}{a hash table always}{a sequence of conditional jumps only}
\end{question}

\begin{question}{M}
Which of the following is NOT a form of three-address instruction?
\optsshort{x = y op z}{goto L}{x = op y}{x = y op z op w}
\end{question}

\begin{question}{M}
The triple representation of a = b * $-$c + b * $-$c needs
\optsshort{6 entries}{5 entries}{7 entries}{4 entries}
\end{question}

\begin{question}{M}
Given below are two statements.\\
Statement I: A syntax tree for an expression always has fewer nodes than the corresponding DAG.\\
Statement II: A DAG for an expression identifies common subexpressions.
\optslong{Both Statement I and Statement II are true}{Both Statement I and Statement II are false}{Statement I is true but Statement II is false}{Statement I is false but Statement II is true}
\end{question}

\begin{question}{L}
Three-address code is called so because each instruction has
\optslong{three temporaries}{exactly three operators}{at most three addresses (two operands and one result)}{three labels}
\end{question}

\begin{question}{M}
In translating the statement while (a < b) a = a + 1;, the intermediate code needs
\optslong{a call instruction}{only one label}{no jumps}{a conditional jump out of the loop and an unconditional jump back to the test}
\end{question}

\begin{answers}
\ans{1}{1}{t1 = minus c; t2 = b * t1; t3 = minus c; t4 = b * t3; t5 = t2 + t4; a = t5.}
\ans{2}{1}{t1 = minus c; t2 = b * t1; t3 = t2 + t2; a = t3.}
\ans{3}{2}{Common implementations of three-address code.}
\ans{4}{1}{Indirect triples avoid this by reordering a list of pointers.}
\ans{5}{2}{The leaves a, b, c, d appear once each, and b $-$ c is shared: 4 leaves + 5 interior nodes.}
\ans{6}{1}{t1 = a + b; t2 = c + d; t1 = t1 * t2.}
\ans{7}{2}{Each assignment defines a new version of the variable.}
\ans{8}{4}{x3 = $\varphi$(x1, x2) selects the value according to the path taken.}
\ans{9}{2}{Each row has 20 elements; the compiler generates this computation in the intermediate code.}
\ans{10}{4}{Lists of jumps (truelist, falselist, nextlist) are patched once the target labels are known.}
\ans{11}{3}{This is the semantics of || in C and "or else" in other languages.}
\ans{12}{1}{Syntax trees, DAGs, postfix and three-address code are all intermediate representations.}
\ans{13}{3}{* has higher precedence than +.}
\ans{14}{2}{The second argument of call is the number of parameters.}
\ans{15}{1}{The same optimizer can be reused for many front ends and back ends.}
\ans{16}{2}{For sparse values, a sequence of tests or a hash table is used instead.}
\ans{17}{4}{Three-address code has at most three addresses and one operator per instruction.}
\ans{18}{1}{(0) minus c; (1) * b, (0); (2) minus c; (3) * b, (2); (4) + (1), (3); (5) = a, (4).}
\ans{19}{4}{A DAG never has more nodes than the syntax tree, because it shares repeated subtrees.}
\ans{20}{3}{For example t1 = b * c.}
\ans{21}{4}{L1: if a >= b goto L2; a = a + 1; goto L1; L2:.}
\end{answers}

\section{Code Optimization and Code Generation (basic blocks, flow graphs, data-flow analysis, local, global, loop and peephole optimization, instruction scheduling)}

\begin{question}{H}
Consider the three-address code: (1) a = 0  (2) i = 1  (3) t = a + i  (4) a = t  (5) i = i + 1  (6) if i <= 10 goto (3)  (7) print a. The number of basic blocks is
\optsshort{2}{4}{7}{3}
\end{question}

\begin{question}{H}
Which of the following instructions are leaders when partitioning three-address code into basic blocks?\\
A. The first instruction  B. Any target of a conditional or unconditional jump  C. Any instruction immediately following a jump  D. Every assignment instruction\\
Choose the correct answer from the options given below:
\optsshort{B, C and D only}{A and D only}{A, B, C and D}{A, B and C only}
\end{question}

\begin{question}{H}
For the basic block a = b + c; b = a $-$ d; c = b + c; d = a $-$ d, the DAG has how many nodes?
\optsshort{8}{6}{7}{5}
\end{question}

\begin{question}{H}
In the basic block a = b + c; b = a $-$ d; c = b + c; d = a $-$ d, which optimization is possible?
\optslong{The statement d = a $-$ d can be replaced by d = b, since a $-$ d was already computed into b}{Loop unrolling}{Removing the first statement}{Constant folding of b + c}
\end{question}

\begin{question}{H}
For the block (1) a = b + c  (2) d = a $-$ e  (3) b = d + f  (4) return b, the set of variables live at the entry of the block is
\optsshort{\{b, c, d, e, f\}}{\{a, d, e, f\}}{\{b, c, e, f\}}{\{a, b, c\}}
\end{question}

\begin{question}{H}
Match List I with List II.\\
List I (Data-flow problem): A. Reaching definitions  B. Live variables  C. Available expressions  D. Very busy expressions\\
List II (Direction and meet operator): I. Backward, intersection  II. Forward, intersection  III. Backward, union  IV. Forward, union
\optsshort{A-III, B-IV, C-II, D-I}{A-IV, B-III, C-I, D-II}{A-IV, B-II, C-III, D-I}{A-IV, B-III, C-II, D-I}
\end{question}

\begin{question}{H}
Match List I with List II.\\
List I (Optimization): A. Constant folding  B. Copy propagation  C. Dead-code elimination  D. Strength reduction\\
List II (Example): I. Replacing i $\times$ 4 inside a loop by repeated addition of 4  II. Replacing x = 2 $\times$ 3.14 by x = 6.28  III. Removing an assignment whose result is never used  IV. After y = x, replacing later uses of y by x
\optsshort{A-IV, B-II, C-III, D-I}{A-II, B-III, C-IV, D-I}{A-II, B-IV, C-III, D-I}{A-II, B-IV, C-I, D-III}
\end{question}

\begin{question}{H}
Moving a computation whose value does not change inside a loop to just before the loop is called
\optslong{loop fusion}{code motion (loop-invariant code motion)}{strength reduction}{loop unrolling}
\end{question}

\begin{question}{H}
Given below are two statements.\\
Statement I: Induction-variable elimination increases the number of variables used in a loop.\\
Statement II: Loop unrolling increases the number of iterations of a loop.
\optslong{Both Statement I and Statement II are true}{Both Statement I and Statement II are false}{Statement I is true but Statement II is false}{Statement I is false but Statement II is true}
\end{question}

\begin{question}{H}
Match List I with List II.\\
List I (Loop optimization): A. Loop unrolling  B. Loop fusion (jamming)  C. Loop-invariant code motion  D. Loop fission\\
List II (Description): I. Moving computations that do not change inside the loop to outside it  II. Splitting one loop into several loops over the same range  III. Replicating the loop body to reduce loop-control overhead  IV. Combining adjacent loops over the same range into one
\optsshort{A-III, B-I, C-IV, D-II}{A-III, B-IV, C-I, D-II}{A-III, B-IV, C-II, D-I}{A-IV, B-III, C-I, D-II}
\end{question}

\begin{question}{H}
Peephole optimization
\optslong{is done only on the source program}{examines a small window of target (or intermediate) instructions and replaces them with shorter or faster sequences}{builds the symbol table}{examines the whole program at once}
\end{question}

\begin{question}{H}
Given below are two statements: one is labelled as Assertion A and the other is labelled as Reason R.\\
Assertion A: The instruction sequence MOV R0, a followed immediately by MOV a, R0 can be reduced to the first instruction alone.\\
Reason R: The second instruction loads a value different from the one just stored.
\optslong{Both A and R are true and R is the correct explanation of A}{Both A and R are true but R is NOT the correct explanation of A}{A is true but R is false}{A is false but R is true}
\end{question}

\begin{question}{M}
Replacing a jump to an unconditional jump (goto L1 where L1: goto L2) by goto L2 is an example of
\optsshort{flow-of-control optimization}{constant folding}{register allocation}{algebraic simplification}
\end{question}

\begin{question}{M}
Register allocation is often modelled as
\optslong{a network-flow problem}{sorting}{a shortest-path problem}{graph colouring of the interference graph, with k colours for k registers}
\end{question}

\begin{question}{M}
Two variables interfere in the interference graph if
\optslong{they are declared in the same block}{they have the same type}{they are live at the same time}{they are never used}
\end{question}

\begin{question}{M}
Given below are two statements.\\
Statement I: Local optimization works within a single basic block.\\
Statement II: Global optimization works across basic blocks, using data-flow analysis on the flow graph.
\optslong{Both Statement I and Statement II are true}{Both Statement I and Statement II are false}{Statement I is true but Statement II is false}{Statement I is false but Statement II is true}
\end{question}

\begin{question}{M}
Instruction scheduling is performed to
\optslong{build parse trees}{allocate memory}{reorder instructions to avoid pipeline stalls and exploit instruction-level parallelism, while preserving dependences}{reduce program size}
\end{question}

\begin{question}{M}
Which of the following is a machine-dependent optimization?
\optslong{Common subexpression elimination}{Peephole optimization using a special increment instruction}{Dead-code elimination}{Constant folding}
\end{question}

\begin{question}{M}
In the flow graph of the code (1) a = 0  (2) i = 1  (3) t = a + i  (4) a = t  (5) i = i + 1  (6) if i <= 10 goto (3)  (7) print a, the number of edges between basic blocks is
\optsshort{2}{4}{5}{3}
\end{question}

\begin{question}{M}
A natural loop in a flow graph is identified by
\optslong{an edge to the entry node}{a back edge n $\rightarrow$ d, where d dominates n}{a basic block with one instruction}{any cycle}
\end{question}

\begin{question}{M}
Given below are two statements: one is labelled as Assertion A and the other is labelled as Reason R.\\
Assertion A: Dead-code elimination can remove the assignment x = y + z if x is not live after it.\\
Reason R: Every variable is live at the end of a basic block.
\optslong{Both A and R are true and R is the correct explanation of A}{Both A and R are true but R is NOT the correct explanation of A}{A is true but R is false}{A is false but R is true}
\end{question}

\begin{question}{M}
Algebraic simplification can replace
\optslong{x = y by x = y + 0}{constants by variables}{loops by recursion}{x = x + 0 and x = x $\times$ 1 by nothing, and x $\times$ 2 by x + x or a shift}
\end{question}

\begin{question}{M}
Next-use information in code generation is used to
\optslong{decide which values to keep in registers and when a register can be freed}{parse the program}{detect lexical errors}{build the symbol table}
\end{question}

\begin{question}{M}
Given below are two statements.\\
Statement I: Using Ershov numbers with every leaf labelled 1, the expression tree for (a + b) $\times$ (c + d) needs 3 registers to be evaluated without storing intermediate values in memory.\\
Statement II: Evaluating the subtree that needs more registers first can reduce the number of registers needed.
\optslong{Both Statement I and Statement II are true}{Both Statement I and Statement II are false}{Statement I is true but Statement II is false}{Statement I is false but Statement II is true}
\end{question}

\begin{question}{M}
Which of the following problems are faced by the code generator?\\
A. Instruction selection  B. Register allocation  C. Evaluation order  D. Lexical analysis\\
Choose the correct answer from the options given below:
\optsshort{A, B and C only}{A, B, C and D}{B, C and D only}{A and D only}
\end{question}

\begin{question}{M}
Common subexpression elimination across basic blocks requires knowing which expressions are
\optsshort{live}{available}{very busy}{dead}
\end{question}

\begin{question}{L}
Unreachable code is code that
\optslong{has no variables}{can never be executed on any path}{is executed most often}{is inside loops}
\end{question}

\begin{question}{M}
Given below are two statements.\\
Statement I: An optimization must preserve the meaning (observable behaviour) of the program.\\
Statement II: Every optimization always makes a program both faster and smaller.
\optslong{Both Statement I and Statement II are true}{Both Statement I and Statement II are false}{Statement I is true but Statement II is false}{Statement I is false but Statement II is true}
\end{question}

\begin{answers}
\ans{1}{4}{Leaders are (1) (first instruction), (3) (a jump target) and (7) (after a jump): blocks \{1, 2\}, \{3--6\} and \{7\}.}
\ans{2}{4}{Each basic block runs from a leader up to (but not including) the next leader.}
\ans{3}{2}{Leaves b0, c0, d0 and interior nodes + (a), $-$ (b and d share it) and + (c): 3 + 3 = 6.}
\ans{4}{1}{This is common subexpression elimination; d has not changed between the two computations of a $-$ d.}
\ans{5}{3}{a and d are defined before they are used; b is used in (1) before it is redefined in (3).}
\ans{6}{4}{"May" problems use union; "must" problems use intersection.}
\ans{7}{3}{Common machine-independent optimizations.}
\ans{8}{2}{The computation is placed in a preheader block.}
\ans{9}{2}{Induction-variable elimination removes redundant variables, and unrolling reduces the number of iterations.}
\ans{10}{2}{Loop transformations improve speed and locality.}
\ans{11}{2}{Examples: removing redundant loads and stores, unreachable code, jumps to jumps, algebraic simplifications and using machine idioms.}
\ans{12}{3}{Provided the second instruction is not a jump target, it reloads exactly the value already in R0.}
\ans{13}{1}{It may also make L1 unreachable, enabling dead-code removal.}
\ans{14}{4}{Optimal graph colouring is NP-complete, so heuristics are used, and spilling handles failures.}
\ans{15}{3}{Interfering variables cannot share a register.}
\ans{16}{1}{Loop optimizations are a major part of global optimization.}
\ans{17}{3}{List scheduling on the data-dependence graph is a common technique.}
\ans{18}{2}{Using machine idioms such as INC depends on the target instruction set.}
\ans{19}{4}{B1 $\rightarrow$ B2, B2 $\rightarrow$ B2 (the loop back edge) and B2 $\rightarrow$ B3.}
\ans{20}{2}{The loop consists of d plus all nodes that can reach n without passing through d.}
\ans{21}{3}{R is false; liveness analysis computes which variables are live, which is what makes A safe.}
\ans{22}{4}{These identities are applied during local optimization or peephole optimization.}
\ans{23}{1}{It is computed by a backward scan over each basic block.}
\ans{24}{1}{Each of a + b and c + d is labelled 2; equal labels give the parent 2 + 1 = 3.}
\ans{25}{1}{Lexical analysis is a front-end phase.}
\ans{26}{2}{An expression is available at a point if it has been computed on every path to it without later changes to its operands.}
\ans{27}{2}{For example, statements after an unconditional return or jump with no label.}
\ans{28}{3}{Some optimizations, such as loop unrolling and inlining, trade size for speed.}
\end{answers}
